% !TeX program = pdflatex
%
% "The fifteen sporadic Apery-like pairs: second solutions, the chi-twisted
%  Lucas law, and the closed forms" -- status paper of River's odd-zeta programme.
%
% Provenance of every statement: work/LBW_GENERAL.md (T1-T4), work/LEMMAD_CHECK.md,
% work/SPORADIC_BARE.md, work/MINIMAL_FORM_PROOF.md (7-9), lean/ZetaLucas.
% Companion papers: papers_out/lucas2nd, papers_out/lucas_min.
%
% Evidence labels are load-bearing and are never silently upgraded; see 1.4.
%
\documentclass[11pt,reqno]{amsart}

\usepackage[T1]{fontenc}
\usepackage[utf8]{inputenc}
\usepackage{amsmath,amssymb,amsthm}
\usepackage{mathtools}
\usepackage{booktabs}
\usepackage{array}
\usepackage{longtable}
\usepackage{enumitem}
\usepackage[margin=1.05in]{geometry}
\usepackage{microtype}
\usepackage[hidelinks]{hyperref}

\theoremstyle{plain}
\newtheorem{theorem}{Theorem}[section]
\newtheorem{proposition}[theorem]{Proposition}
\newtheorem{lemma}[theorem]{Lemma}
\newtheorem{corollary}[theorem]{Corollary}
\newtheorem{conjecture}[theorem]{Conjecture}

\theoremstyle{definition}
\newtheorem{definition}[theorem]{Definition}
\newtheorem{problem}[theorem]{Problem}

\theoremstyle{remark}
\newtheorem{remark}[theorem]{Remark}
\newtheorem{observation}[theorem]{Observation}

\newcommand{\Z}{\mathbb{Z}}
\newcommand{\Q}{\mathbb{Q}}
\newcommand{\R}{\mathbb{R}}
\newcommand{\Zp}{\Z_{(p)}}
\newcommand{\vp}{v_p}
\newcommand{\bin}[2]{\binom{#1}{#2}}
\newcommand{\HH}[2]{H^{(#1)}_{#2}}
\newcommand{\KK}[3]{K^{(#1)}_{#2}\!\left(#3\right)}
\DeclareMathOperator{\lcm}{lcm}
\DeclareMathOperator{\ct}{ct}

% evidence-class markers
\newcommand{\Thm}{\textup{\textsf{[THM]}}}
\newcommand{\Proved}{\textup{\textsf{[PROVED]}}}
\newcommand{\Lean}{\textup{\textsf{[LEAN]}}}
\newcommand{\Verified}{\textup{\textsf{[VERIFIED]}}}
\newcommand{\VerifiedR}[1]{\textup{\textsf{[VERIFIED: #1]}}}
\newcommand{\Excluded}{\textup{\textsf{[EXCLUDED]}}}
\newcommand{\Open}{\textup{\textsf{[OPEN]}}}

\numberwithin{equation}{section}

\begin{document}

\title[The fifteen sporadic Ap\'ery-like pairs]
      {The fifteen sporadic Ap\'ery-like pairs:\\
       second solutions, the $\chi$-twisted Lucas law,\\ and the closed forms}

\author{}
\date{\today}

\subjclass[2020]{11A07, 11B65, 11S80, 33F10, 05A10}
\keywords{Ap\'ery-like sequences, sporadic sequences, Lucas congruences, second
solution, Ap\'ery limits, Dwork crystals, creative telescoping, Gosper's algorithm}

\begin{abstract}
For each of the fifteen sporadic Ap\'ery-like sequences $A(n)$ --- Zagier's six, the six
Almkvist--Zudilin cases and Cooper's three --- let $B(n)$ be the second solution of the
same three-term recurrence, normalised by $B(0)=0$, $B(1)=1$. The first row obeys a Lucas
congruence $A(ap+r)\equiv A(a)A(r)\pmod p$. We report the current state of the second row.
Exact computation across all fifteen pairs and all primes $5\le p\le 103$ says that the
naive analogue $p^{w}B(ap+r)\equiv B(a)A(r)$ is \emph{false} for exactly seven of the
fifteen, at exactly the primes at which a quadratic Dirichlet character $\chi$ takes the
value $-1$, and that inserting the single factor $\chi(p)$ repairs it with $p$-adic floor
exactly $w$ in every cell tested. We isolate the mechanism (one line of Frobenius descent
for a $\chi$-twisted harmonic letter), prove the descent under five explicit hypotheses,
and correct the product-region hypothesis of the earlier statement: the surviving digit set
is a proper subset and carries a base-level defect $\Delta(a)$, whose vanishing is an
additional clause. The clause is automatic for tame weights, fails for $\alpha$
(Domb), $\varepsilon$ and $\mathbf E$, and \emph{no admissible weight repairs it} --- for
$\mathbf E$, the only known instance with $\chi\ne1$, the defect is independent of the
choice of decomposition. Four instances are theorems ($\gamma$, $\mathbf A$, $\mathbf D$,
$s_{10}$), one ($s_7$) is one bounded local lemma away, three are blocked as above, and
seven remain conjectural. Separately, a step-three endpoint transform proves the previously computational
denominator bound $d_n^2B(n)\in\Z$ for Zagier's $\mathbf B$, shows the exponent two is
optimal, and classifies its full dilation family by the exact condition $3\mid h$.
On the closed-form side we record a proof, formalised in Lean~4,
of $b_n=\sum_k\bin nk^2\bin{n+k}k^2(2\HH3n-\HH3k)$ for Ap\'ery's $\zeta(3)$ numerators;
finite-certificate proofs, valid for every $n$, of the symmetric closed forms for the Franel
numbers and for $\sum_k\bin nk^4$ --- these sums are the $t^2$-coefficients of
Straub--Zudilin's deformation, and we say so, together with an elementary re-derivation of
the Ap\'ery limits $\zeta(2)/4$ and $\zeta(2)/5$ that they proved for all exponents; a proof
that the mechanism behind those two \emph{cannot} close the $\zeta(2)$-Ap\'ery case; and a
$\chi_{-4}$-twisted closed form for the Catalan-constant family $\mathbf E$. Three of the fifteen have no archimedean
Ap\'ery limit at all and satisfy the $p$-adic law regardless: the descent is a statement
about Frobenius. We locate the law precisely against Beukers--Vlasenko's Conjecture 7.5:
the two are the two rows of one $2\times2$ Frobenius matrix, neither implies the other, and
the character is invisible in unit-root formulations.
\end{abstract}

\maketitle

%======================================================================
\section{Introduction}

\subsection{The two rows}
Ap\'ery's proof that $\zeta(3)\notin\Q$ rests on a pair of solutions of one three-term
recurrence: the integral solution $a_n=\sum_k\bin nk^2\bin{n+k}k^2$ and a companion $b_n$
with $b_n/a_n\to\zeta(3)$. The arithmetic of the two rows is not symmetric. The first row
is integral and satisfies a Lucas congruence,
\begin{equation}\label{eq:Arow}
  A(ap+r)\;\equiv\;A(a)A(r)\pmod p\qquad(0\le a,r<p),
\end{equation}
proved for $a_n$ by Gessel \cite{Gessel1982}, extended to the sporadic class by
Malik--Straub \cite{MalikStraub}, obtained for fourteen of the fifteen in a strong form by
Gorodetsky \cite{Gorodetsky}, and completed for all fifteen by Straub \cite{Straub2301},
who resolves the two remaining cases with constant-term representations. The second row is not integral --- it has denominators
$d_n^{\,w}$, $d_n=\lcm(1,\dots,n)$ --- and no congruence of Lucas type for it appears in
the literature.

This paper is a status report on the second row for the whole sporadic class. It collects
what is proved, states honestly what is only computed, and --- this is its main technical
purpose --- records two corrections to the companion papers
\cite{CompanionLucas2nd,CompanionLucasMin}, and assembles the closed forms that the
congruence machinery consumes.

\subsection{The law}
Write $q=\lfloor n/p\rfloor$. For each sporadic pair there is an integer $w\in\{2,3\}$ and
a quadratic Dirichlet character $\chi$ (possibly trivial) such that, empirically,
\begin{equation}\label{eq:master}
  \vp\bigl(p^{w}B(n)A(q)-\chi(p)B(q)A(n)\bigr)\;\ge\;w ,
  \tag{$\mathrm{LB}_w^\chi$}
\end{equation}
with equality in every cell tested. The untwisted form ($\chi\equiv1$) is
\emph{false} for seven of the fifteen families, and the failing primes are exactly those
with $\chi(p)=-1$; there the valuation drops to $0$ in the single-digit range and as low as
the integrality bound permits beyond it. The character is the character of the sequence's
Ap\'ery limit: the limit is a rational multiple of $\zeta(w)$ when $\chi$ is trivial and of
$L(\chi,w)$ otherwise.

We prove the single-digit, mod-$p$ shadow of \eqref{eq:master},
\begin{equation}\label{eq:LB}
  p^{w}B(ap+r)\;\equiv\;\chi(p)^{e}\,B(a)\,A(r)\pmod p\qquad(1\le a<p,\ 0\le r<p),
\end{equation}
under hypotheses (H1)--(H5) of \S\ref{sec:LB}, and verify those hypotheses for four of the
fifteen. The full statement \eqref{eq:master} --- all $n$, modulus $p^w$ --- is
\Verified\ only and is stated as Conjecture~\ref{conj:master}; we never call it a theorem.

\subsection{What is new here}
\begin{enumerate}[leftmargin=2em]
\item \textbf{A correction to the product-region hypothesis} (\S\ref{sec:LB}). The
  hypothesis (H2) used in \cite{CompanionLucas2nd} asserts that the set of surviving
  base-$p$ digits is $\{0\le b\le a\}\times\Sigma_r$. That is false: the $b$-side is a
  proper subset for every $a\ge(p+1)/2$, already for the Ap\'ery $\zeta(3)$ pair, and the
  omitted terms contribute a \emph{base-level defect} $\Delta(a)$ whose vanishing mod $p$
  is a further clause. The clause is automatic under tameness, so the four proved instances
  are unaffected; it is exactly what fails for $\alpha$, $\varepsilon$, $\mathbf E$.
\item \textbf{A proved negative} (\S\ref{sec:blocked}). No admissible harmonic weight
  repairs $\Delta(a)$ for $\alpha$, $\varepsilon$ or $\mathbf E$; for $\mathbf E$ the class
  $\Delta(a)\bmod p$ is an invariant of the family, not of the fitted weight. Consequently
  the earlier plan --- ``prove one valuation inequality and move four families inside the
  theorem'' --- is wrong; it works for $s_7$ alone.
\item \textbf{Two closed forms with complete finite certificates} (\S\ref{sec:closed}).
  Symmetric harmonic weights for the Franel numbers and for $\sum_k\bin nk^4$, with every
  certificate identity re-verified symbolically here, valid for every $n\ge1$. These sums
  are \emph{not} new: they are the $t^2$-coefficients of Straub--Zudilin's analytic
  deformation \cite{StraubZudilin}, as we show in Remark~\ref{rem:SZ}, and their Ap\'ery
  limits $\zeta(2)/(d+1)$ are theorems of \emph{loc.\ cit.} What is contributed here is the
  finite certificate proof --- Straub--Zudilin obtain the recurrence only for large enough
  $n$, and their route passes through a locally uniform limit --- and the explicit
  harmonic-monomial display, which is what Theorem~\ref{thm:LB} consumes.
\item \textbf{A provable insufficiency} (\S\ref{sec:Dobstruction}). The mechanism that
  proves those two closed forms provably fails for the $\zeta(2)$-Ap\'ery pair: one
  component of the correction is not Gosper-summable. This is the first place in the
  programme where the rank-two mechanism is \emph{provably} insufficient.
\end{enumerate}

\subsection{Evidence classes}\label{sec:labels}
The programme's labelling discipline is retained verbatim in this paper, and the labels are
load-bearing. \Thm: proved and adversarially refereed. \Proved: a complete proof is given
or cited to a written proof. \Lean: checked by the Lean~4 kernel with no \texttt{sorry}.
\Verified\ (with a range): an exact finite computation, \emph{never} treated as proof.
\Excluded: a proved negative. \Open. A statement carrying \Verified\ is never restated
later as though it carried \Proved.

\subsection{Relation to the companion papers}
The weight-three congruence \eqref{eq:LB} for the Ap\'ery $\zeta(3)$ pair, its elementary
self-contained proof, the sharpness analysis at $p=2,3$ and the general theorem are the
subject of \cite{CompanionLucas2nd}; the closed form for $b_n$ together with its Lean
formalisation is \cite{CompanionLucasMin}. We do not repeat those proofs. What is here is
the class-wide picture, the corrections of items~1--2 above, and the closed-form material
of \S\ref{sec:closed}.

\subsection{Conventions}
$p$ is a prime, $\Zp$ the localisation of $\Z$ at $p$, $\vp$ the $p$-adic valuation on
$\Q$, and $x\equiv y\pmod{p^m}$ for $x,y\in\Q$ means $\vp(x-y)\ge m$. We write
$\HH ry=\sum_{m=1}^y m^{-r}$ and, for a Dirichlet character $\chi$,
$\KK r\chi y=\sum_{m=1}^{y}\chi(m)m^{-r}$; thus $\HH ry=\KK r{\mathbf 1}y$. $\chi_D$ is the
Kronecker symbol $\bigl(\tfrac D\cdot\bigr)$, and $L(\chi,s)=\sum_{m\ge1}\chi(m)m^{-s}$.

%======================================================================
\section{The fifteen pairs}\label{sec:pairs}

\subsection{The recurrences}
Two families of three-term recurrences, with $u_{-1}=0$, $u_0=1$:
\begin{align}
  (n+1)^2u_{n+1}&=(an^2+an+b)u_n-cn^2u_{n-1},\label{eq:R2}\\
  (n+1)^3u_{n+1}&=(2n+1)(an^2+an+b)u_n-n(cn^2+d)u_{n-1}.\label{eq:R3}
\end{align}
Up to degenerate cases, \eqref{eq:R2} has six sporadic integral solutions, found by Zagier
\cite{Zagier2009} and labelled $\mathbf A,\dots,\mathbf F$; \eqref{eq:R3} with $d=0$ has six,
found by Almkvist--Zudilin \cite{AlmkvistZudilin2006} and labelled
$\alpha,\gamma,\delta,\varepsilon,\zeta,\eta$; and Cooper \cite{Cooper2012} found three
further solutions with $d\ne0$ (so counted in \cite[\S2]{MalikStraub}; \cite{Cooper2012}
itself announces two), labelled $s_7$, $s_{10}$, $s_{18}$ after the levels of the
associated modular forms; they satisfy two-term supercongruences \cite{OsburnSahuStraub},
and Cooper has since extended the search to four-term recurrences \cite{Cooper2302}. We
follow the labelling and the binomial sums of Malik--Straub \cite[Tables 1--2]{MalikStraub},
whose second column records the alternative Almkvist--van~Straten--Zudilin labels
$(\mathrm a),\dots,(\mathrm g)$.

\begin{definition}\label{def:B}
For each pair, $B(n)$ is the second solution: $B(0)=0$, $B(1)=1$, and the recurrence is
imposed for $n\ge1$ (the step $n=0$ is degenerate). This is Ap\'ery's own normalisation and
that of Chamberland--Straub \cite{ChamStraub}. Rescaling $B$ by a $p$-unit rational does not
affect any statement below, so for $p\ge5$ the normalisation is harmless.
\end{definition}

\subsection{The table}
Table~\ref{tab:fifteen} lists the fifteen pairs with their parameters, first row, Ap\'ery
limit $L=\lim B(n)/A(n)$, weight $w$ and character $\chi$.

\begin{table}[t]
\small
\setlength{\tabcolsep}{4pt}
\begin{tabular}{@{}rllll@{\ }cc@{}}
\toprule
\# & family & $(a,b,c[,d])$ & $A(n)$ & $L$ & $w$ & $\chi$\\
\midrule
1&$\mathbf A$ (a)&$(7,2,-8)$&$\sum_k\bin nk^3$ (Franel)&$\zeta(2)/4$&2&$\mathbf 1$\\
2&$\mathbf B$ (f)&$(9,3,27)$&$\sum_k(-1)^k3^{n-3k}\bin n{3k}\tfrac{(3k)!}{k!^3}$&none&2&$\chi_{-3}$\\
3&$\mathbf C$ (c)&$(10,3,9)$&$\sum_k\bin nk^2\bin{2k}k$&$L(\chi_{-3},2)/2$&2&$\chi_{-3}$\\
4&$\mathbf D$ (b)&$(11,3,-1)$&$\sum_k\bin nk^2\bin{n+k}n$&$\zeta(2)/5$&2&$\mathbf 1$\\
5&$\mathbf E$ (d)&$(12,4,32)$&$\sum_k\bin nk\bin{2k}k\bin{2(n-k)}{n-k}$&$G/2$&2&$\chi_{-4}$\\
6&$\mathbf F$ (g)&$(17,6,72)$&$\sum_k(-1)^k8^{n-k}\bin nk\sum_l\bin kl^3$&$\tfrac58L(\chi_{-3},2)$&2&$\chi_{-3}$\\
7&$\alpha$&$(10,4,64,0)$&$\sum_k\bin nk^2\bin{2k}k\bin{2(n-k)}{n-k}$ (Domb)&$7\zeta(3)/24$&3&$\mathbf 1$\\
8&$\gamma$&$(17,5,1,0)$&$\sum_k\bin nk^2\bin{n+k}n^2$ (Ap\'ery)&$\zeta(3)/6$&3&$\mathbf 1$\\
9&$\delta$&$(7,3,81,0)$&$\sum_k(-1)^k3^{n-3k}\bin n{3k}\bin{n+k}n\tfrac{(3k)!}{k!^3}$&none&3&$\mathbf 1$\\
10&$\varepsilon$&$(12,4,16,0)$&$\sum_k\bin nk^2\bin{2k}n^2$&$7\zeta(3)/32$&3&$\mathbf 1$\\
11&$\zeta$&$(9,3,-27,0)$&$\sum_{k,l}\bin nk^2\bin nl\bin kl\bin{k+l}n$&$L(\chi_{-3},3)/3$&3&$\chi_{-3}$\\
12&$\eta$&$(11,5,125,0)$&$\sum_k(-1)^k\bin nk^3\bigl[\bin{4n-5k-1}{3n}+\bin{4n-5k}{3n}\bigr]$&none&3&$\chi_{5}$\\
13&$s_7$&$(13,4,-27,3)$&$\sum_k\bin nk^2\bin{n+k}k\bin{2k}n$&$\zeta(2)/7$&2&$\mathbf 1$\\
14&$s_{10}$&$(6,2,-64,4)$&$\sum_k\bin nk^4$&$\zeta(2)/5$&2&$\mathbf 1$\\
15&$s_{18}$&$(14,6,192,-12)$&see \eqref{eq:s18}&$L(\chi_{-3},2)/2$&2&$\chi_{-3}$\\
\bottomrule
\end{tabular}
\medskip
\caption{The fifteen sporadic pairs. Parenthesised letters in column~2 are the
Almkvist--van~Straten--Zudilin labels of \cite[Table~1]{MalikStraub}. $G=L(\chi_{-4},2)$ is
Catalan's constant. Rows 1--6 satisfy \eqref{eq:R2}; rows 7--12 satisfy \eqref{eq:R3} with
$d=0$; rows 13--15 satisfy \eqref{eq:R3} with $d\ne0$.}
\label{tab:fifteen}
\end{table}

Here $s_{18}(0)=1$ and, for $n\ge1$,
\begin{equation}\label{eq:s18}
  s_{18}(n)=\sum_{k=0}^{\lfloor n/3\rfloor}(-1)^k\bin nk\bin{2k}k\bin{2(n-k)}{n-k}
     \Bigl[\bin{2n-3k-1}n+\bin{2n-3k}n\Bigr].
\end{equation}

\VerifiedR{$n\le12$, exact} each binomial sum in Table~\ref{tab:fifteen} satisfies its
recurrence, re-checked independently for this paper.

\subsection{The limits}
\VerifiedR{$400$ digits} The limits in Table~\ref{tab:fifteen} were identified by two-term
integer relation detection against a $33$-constant basis
($\zeta(2),\zeta(3),\zeta(4)$; $L(\chi_D,2)$, $L(\chi_D,3)$ for
$D=-3,-4,-7,-8,-11,-15,-20,-24,5,8,12,13,24$; $\log^22$, $\pi^2\log2$, $\log^32$), matching
to between $31$ and $401$ digits (the slowest, $\mathbf F$, converges at rate $(8/9)^n$).
\emph{We claim no novelty for any of these values.} Almkvist--van~Straten--Zudilin
\cite[\S2]{ASvSZ} tabulate the limits of the six weight-two families --- $\pi^2/24$
($\mathbf A$), $\pi^2/30$ ($\mathbf D$), $L(\chi_{-3},2)/2$ ($\mathbf C$), $G/2$
($\mathbf E$), $\tfrac58L(\chi_{-3},2)$ ($\mathbf F$) --- and record that the remaining
one, $\mathbf B$, is \emph{not convergent}, attributing the proofs to Zagier
\cite{Zagier2009} and Yang \cite{Yang2008}. The Franel and $\sum_k\bin nk^4$ limits are
also \cite[eq.~(20)]{ChamStraub} (Cusick, via \cite[p.~202]{vdPoorten1979}) and are proved
in general by Straub--Zudilin \cite{StraubZudilin}; $\zeta(3)/6$ is Ap\'ery's; and the Domb
limit $7\zeta(3)/24$, for exactly the normalisation of Definition~\ref{def:B}, is proved by
Shvets \cite{Shvets2026}. The values for $\varepsilon$, $\zeta$, $s_7$, $s_{18}$ and the
non-convergence of $\delta$ and $\eta$ we have not located in the literature, but they are
of the shape predicted by \cite{Zagier2009,AlmkvistZudilin2006} and we record them only as
\Verified.

\subsection{The weights, and their sharpness}
\begin{observation}\label{obs:weight}
\VerifiedR{$n\le600$; independently $n\le120$ for this paper}
Let $w$ be the least exponent with $d_n^{\,w}B(n)\in\Z$ for all $n$ in the tested range,
$d_n=\lcm(1,\dots,n)$. Then
\begin{align*}
  w&=2 &&\text{for all six of \eqref{eq:R2} \emph{and} all three of Cooper's $d\ne0$
     cases},\\
  w&=3 &&\text{for all six Almkvist--Zudilin $d=0$ cases}.
\end{align*}
It is sharp: $d_n^{\,w-1}B(n)\notin\Z$ for at least $596$ of the $600$ values of $n$ in the
first sweep (and for $116$--$119$ of the $120$ in the second). Per prime,
$\vp(B(n))\ge-w\lfloor\log_pn\rfloor$, \VerifiedR{$0$ failures, $15$ families $\times$
$16$ primes $\times$ $n\le400$}.
\end{observation}

Two things are worth isolating. First, $w$ tracks the Cooper deformation $d\ne0$, not the
shape of the recurrence: Cooper's three live in the cubic family \eqref{eq:R3} and have
weight $2$, not $3$. Second, wherever the Ap\'ery limit exists it is a rational multiple of
$\zeta(w)$ or $L(\chi,w)$, so $w$ agrees with the motivic weight of the limit --- $12$
cases out of $12$.

\subsection{Three pairs with no archimedean limit}\label{sec:nolimit}
\begin{proposition}\label{prop:nolimit}
\Proved\ For $\mathbf B$, $\delta$ and $\eta$ the characteristic equation of the recurrence
has a pair of complex-conjugate roots of equal modulus, with discriminants
\[
  \mathbf B:\ a^2-4c=-27,\qquad
  \delta:\ 4a^2-4c=-128,\qquad
  \eta:\ 4a^2-4c=-16 ,
\]
so that Poincar\'e--Perron provides no dominant solution and $B(n)/A(n)$ does not converge.
Every other sporadic pair has real distinct characteristic roots.
\end{proposition}

\begin{proof}
The characteristic equation of \eqref{eq:R2} is $\lambda^2-a\lambda+c=0$ and of
\eqref{eq:R3} is $\lambda^2-2a\lambda+c=0$; substitute the parameters of
Table~\ref{tab:fifteen}. The discriminants of the remaining twelve are
$81,64,125,16,1,144,1152,512,432,784,400,16$, all positive.
\end{proof}

\VerifiedR{$n=590$ vs.\ $n=600$} the ratios $B(n)/A(n)$ for these three agree to $0$ digits,
against $380$ digits for the convergent families. Nevertheless all three satisfy
\eqref{eq:master} flat, with $\chi=\chi_{-3}$, $\mathbf 1$, $\chi_5$ respectively. This is
the cleanest evidence in the programme that \eqref{eq:master} is a statement about
\emph{Frobenius}, not about the archimedean limit; we return to it in \S\ref{sec:BV}.

%======================================================================
\section{The $\chi$-twisted law}\label{sec:law}

\subsection{Where the untwisted law dies}
\VerifiedR{$5\le p\le43$, $n\le300$, exact rational arithmetic}
Testing $\vp(p^wB(n)A(q)-B(q)A(n))\ge w$ with $q=\lfloor n/p\rfloor$:

\begin{center}\small
\begin{tabular}{@{}ll@{}}
\toprule
verdict & families\\
\midrule
holds flat at every prime tested &
  $\mathbf A,\mathbf D,\alpha,\gamma,\delta,\varepsilon,s_7,s_{10}$\quad(8)\\
fails at $p\equiv2\pmod 3$ & $\mathbf B,\mathbf C,\mathbf F,\zeta,s_{18}$\quad(5)\\
fails at $p\equiv3\pmod 4$ & $\mathbf E$\quad(1)\\
fails at $p\equiv\pm2\pmod 5$ & $\eta$\quad(1)\\
\bottomrule
\end{tabular}
\end{center}

The failures are total, not marginal: at a failing prime the single-digit floor is $0$ (the
difference is a $p$-unit) and in the multi-digit range the valuation falls to
$-w\cdot(\#\text{digits}-1)$, i.e.\ as low as Observation~\ref{obs:weight} permits. They
are not explained by exceptional digits ($p\mid A(a)$), by the centre residue, or by
$\vp\bin{2n}n$: the failing set is exactly $\{p:\chi(p)=-1\}$ for one quadratic $\chi$. The
cross-checks that isolate the character are \Verified: the passing primes of the
$\chi_{-3}$ group are mixed mod $4$ and mod $5$; those of $\mathbf E$ are mixed mod $3$ and
mod $5$; and those of $\eta$ in the tested range, $11,19,29,31,41$, are mixed mod $3$ and
mod $4$ and are precisely the $p\equiv\pm1\pmod5$.

\subsection{The repair}
\begin{conjecture}[master form]\label{conj:master}
For every sporadic pair, with $w$ and $\chi$ as in Table~\ref{tab:fifteen}, every prime
$p\ge5$ and every $n\ge1$, writing $q=\lfloor n/p\rfloor$,
\[
  \vp\bigl(p^{w}B(n)A(q)-\chi(p)B(q)A(n)\bigr)\;\ge\;w .
\]
\end{conjecture}

\VerifiedR{$0$ failures} all $15$ pairs, all primes $5\le p\le103$, $n\le400$; and
$n\le1200$ for $p=5,7,11$ --- five base-$p$ digits at $p=5$. In every one of the
$15\times16=240$ (family, prime) cells with $5\le p\le61$ the floor is \emph{exactly} $w$:
perfectly flat, with no dependence on the digit pattern, no centre-residue exception and no
exceptional-digit exception. Independently for this paper: all $15$ families at $p=5$ and
$p=7$, all $n<p^3$, floor exactly $w$ in all $30$ cells.

\begin{remark}[the ramified prime]\label{rem:ramified}
For $\eta$ the prime $p=5$ is ramified, $\chi_5(5)=0$, and Conjecture~\ref{conj:master}
degenerates to $\vp(p^3B(n)A(q))\ge3$, i.e.\ to the $5$-integrality of $B(n)$. That is
true: \VerifiedR{$n\le500$} $\min_n v_5(B(n))=0$, the pole simply does not develop at the
ramified prime, and the floor is exactly $3$. This is the only prime $\ge5$ at which any of
the fifteen characters ramifies; the $\chi_{-3}$ and $\chi_{-4}$ families ramify only at
$3$ and $2$.
\end{remark}

\begin{remark}[$p=2,3$]
Excluded, correctly. With the twisted law and $n\le120$ the floor drops below $w$ for
$\mathbf A$ ($p=2$, floor $1$), $\mathbf F$ ($p=2$, $1$), $\alpha$ ($p=2$ and $p=3$, $2$),
$\gamma$ ($p=3$, $2$), $\delta$ ($p=3$, $1$), $\varepsilon$ ($p=2$, $2$). See
\cite[\S8]{CompanionLucas2nd} for the analysis of where $p\ge5$ enters.
\end{remark}

%======================================================================
\section{Lemma K and Theorem LB}\label{sec:LB}

\subsection{The source of the character}
\begin{lemma}[Lemma K]\label{lem:K}
\Proved\ Let $\chi$ be completely multiplicative (a Dirichlet character; $\chi\equiv1$
allowed), let $r\ge1$ and $y\ge0$. Then, as an identity in $\Q$,
\[
  p^{r}\,\KK r\chi y \;=\; \chi(p)\,\KK r\chi{\lfloor y/p\rfloor}\;+\;p^{r}T,
  \qquad T:=\sum_{m\le y,\ p\nmid m}\frac{\chi(m)}{m^{r}}\in\Zp .
\]
\end{lemma}

\begin{proof}
Split the sum according to $p\mid m$. The $p$-divisible part is
$\sum_{j\le\lfloor y/p\rfloor}\chi(jp)(jp)^{-r}
 =\chi(p)p^{-r}\sum_{j\le\lfloor y/p\rfloor}\chi(j)j^{-r}$, by complete multiplicativity.
The complementary part has $p$-integral terms.
\end{proof}

This one line is the entire source of the factor $\chi(p)$ in \eqref{eq:LB}: a second
solution whose weight needs $\chi$-twisted harmonic letters inherits $\chi(p)$ when the
letters descend, and one whose weight is a pure $\zeta$-harmonic monomial inherits $+1$.

\begin{corollary}\label{cor:K2}
\Proved\ Let $M=\prod_{t=1}^{s}\KK{r_t}{\chi_t}{x_t}$ with $\sum_t r_t=w$, and suppose every
$\KK{r_t}{\chi_t}{\lfloor x_t/p\rfloor}$ lies in $\Zp$. Then
$p^{w}M\equiv\bigl(\prod_t\chi_t(p)\bigr)\prod_t\KK{r_t}{\chi_t}{\lfloor x_t/p\rfloor}
 \pmod p$.
\end{corollary}

\begin{proof}
$p^wM=\prod_t\bigl(\chi_t(p)\KK{r_t}{\chi_t}{\lfloor x_t/p\rfloor}+p^{r_t}T_t\bigr)$ by
Lemma~\ref{lem:K}; expand. Every factor is $p$-integral and every cross term carries a
factor $p$.
\end{proof}

\subsection{The hypotheses, corrected}
Fix a prime $p$ and write $n=ap+r$, $k=bp+s$ with $1\le a<p$, $0\le r,s<p$. Let
\[
  S(n,k)=\prod_i\bin{L_i(n,k)}{M_i(n,k)},\quad A(n)=\sum_kS(n,k),\quad
  B(n)=\sum_kS(n,k)\,\mathfrak w(n,k),
\]
with $L_i,M_i$ integer linear forms, and let the weight be a $\Q$-combination of monomials
of total degree $w$,
\begin{equation}\label{eq:weightshape}
  \mathfrak w(n,k)=\sum_j c_j\prod_t \KK{r_{jt}}{\chi}{x_{jt}(n,k)},\qquad
  \sum_t r_{jt}=w ,
\end{equation}
in letters attached to one fixed quadratic character $\chi$ and to the trivial character,
with $x_{jt}$ integer linear forms. Put $\mathcal B_a=\{b: p\nmid S(a,b)\}$ and define the
\emph{base-level defect}
\begin{equation}\label{eq:Delta}
  \Delta(a)\;:=\;\sum_{b\notin\mathcal B_a}S(a,b)\,\mathfrak w(a,b)
  \;=\;B(a)-\sum_{b\in\mathcal B_a}S(a,b)\,\mathfrak w(a,b).
\end{equation}

\begin{itemize}[leftmargin=2.4em]
\item[(H1)] \emph{Lucas/carry dichotomy.} For every $k$, either $p\mid S(n,k)$ or
  $S(n,k)\equiv S(a,b)S(r,s)\pmod p$.
\item[(H2)] \emph{Product region and base-level defect.} The surviving set
  $\{k:p\nmid S(n,k)\}$ is $\mathcal B_a\times\Sigma_r$ for some $\Sigma_r\subseteq[0,p)$;
  $\sum_{s\in\Sigma_r}S(r,s)\equiv A(r)\pmod p$; and $\Delta(a)\equiv0\pmod p$.
\item[(H3)] \emph{Digit compatibility.} On the surviving set,
  $\lfloor x(n,k)/p\rfloor=x(a,b)$ for every argument form $x$ occurring in
  $\mathfrak w$.
\item[(H4)] \emph{Tameness.} $0\le x(n,k)\le n$ on the summation range for every argument
  form $x$; and $c_j\in\Zp$ for every $j$.
\item[(H5)] \emph{$\chi$-homogeneity.} Every monomial of $\mathfrak w$ contains the same
  number $e\in\{0,1\}$ of $\chi$-letters, the others being trivial-character letters.
\end{itemize}

\begin{remark}[what was corrected, and why it matters]\label{rem:H2}
The version of (H2) used in \cite{CompanionLucas2nd} reads ``the surviving set is
$\{0\le b\le a\}\times\Sigma_r$'', with no $\Delta$-clause. The $b$-side is \emph{not}
$\{0,\dots,a\}$ in general. Already for $\gamma$, with
$S(n,k)=\bin nk^2\bin{n+k}n^2$, one has $p\mid\bin{a+b}a$ as soon as $a+b\ge p$, so digits
$b$ with $b>p-1-a$ drop out, and the $b$-side is proper for every $a\ge(p+1)/2$; the same
happens for $\mathbf D$ and, with a different mechanism, for $\alpha$ (at $p=5$, $a=3$ the
digit $b=0$ is dropped because $5\mid\bin63$). \VerifiedR{all four non-tame families,
$5\le p\le23$} The omitted terms are collected by $\Delta(a)$, and the proof below needs
$\Delta(a)\equiv0$. Lemma~\ref{lem:Deltatame} shows the clause is free under (H4), so the
conclusions of \cite{CompanionLucas2nd} for the four tame families stand as stated; the
clause is nevertheless the exact obstruction for $\alpha$, $\varepsilon$ and $\mathbf E$
(\S\ref{sec:blocked}).
\end{remark}

\begin{lemma}\label{lem:Deltatame}
\Proved\ If (H4) holds then $\mathfrak w(a,b)\in\Zp$ for all $b$, and hence
$\Delta(a)\equiv0\pmod p$ automatically.
\end{lemma}

\begin{proof}
By (H4), every argument satisfies $x(a,b)\le a<p$, so no denominator of
$\KK{r}{\chi}{x(a,b)}$ is divisible by $p$; with $c_j\in\Zp$ this gives
$\mathfrak w(a,b)\in\Zp$. Each term of \eqref{eq:Delta} then has $p\mid S(a,b)$ times a
$p$-integral factor.
\end{proof}

\subsection{The theorem}
\begin{theorem}[Theorem LB]\label{thm:LB}
\Proved\ Under \textup{(H1)--(H5)}, for all $1\le a<p$ and $0\le r<p$,
\[
  p^{w}\,B(ap+r)\;\equiv\;\chi(p)^{e}\,B(a)\,A(r)\pmod p .
\]
\end{theorem}

\begin{proof}
Let $n=ap+r$. By (H4) every argument satisfies $x(n,k)\le n<p^2$, so
$\lfloor x(n,k)/p\rfloor\le a<p$ and each descended letter
$\KK{r}{\chi}{\lfloor x(n,k)/p\rfloor}$ lies in $\Zp$; likewise $x(a,b)\le a<p$, so
$\mathfrak w(a,b)\in\Zp$, and $p^w\mathfrak w(n,k)\in\Zp$ for \emph{every} $k$.
Corollary~\ref{cor:K2} and (H5) give, for every $k$,
\[
  p^{w}\mathfrak w(n,k)\;\equiv\;\chi(p)^{e}\,\widetilde{\mathfrak w}(n,k)\pmod p,
  \qquad
  \widetilde{\mathfrak w}(n,k)
   =\sum_jc_j\prod_t\KK{r_{jt}}{\chi}{\lfloor x_{jt}(n,k)/p\rfloor},
\]
and (H3) upgrades $\widetilde{\mathfrak w}(n,k)=\mathfrak w(a,b)$ on the surviving set. Now
split
\[
  p^{w}B(n)=\sum_{k:\,p\mid S(n,k)}S(n,k)\,p^{w}\mathfrak w(n,k)
           +\sum_{k\ \mathrm{surviving}}S(n,k)\,p^{w}\mathfrak w(n,k).
\]
Every $p^{w}\mathfrak w(n,k)$ is $p$-integral, so the first sum is $\equiv0\pmod p$. In the
second, substitute (H1) and the displayed congruence and use the product structure of (H2):
\[
  p^{w}B(n)\;\equiv\;\chi(p)^{e}\Bigl(\sum_{b\in\mathcal B_a}S(a,b)\mathfrak w(a,b)\Bigr)
                     \Bigl(\sum_{s\in\Sigma_r}S(r,s)\Bigr)
   \;\equiv\;\chi(p)^{e}\bigl(B(a)-\Delta(a)\bigr)A(r)\pmod p ,
\]
and the clause $\Delta(a)\equiv0$ of (H2) finishes.
\end{proof}

\begin{remark}[the hypotheses in context]
(H1) and the product structure in (H2) are exactly the Lucas-factorisation-plus-carry-\-
annihilation property established for all fifteen summands by Malik--Straub
\cite{MalikStraub}, so they are available class-wide. What is \emph{not} needed is a
per-carry valuation bound: (H4) makes $p^{w}\mathfrak w$ $p$-integral outright. Where
tameness fails one needs a substitute, and \S\ref{sec:blocked} shows what that substitute
has to be.
\end{remark}

\begin{remark}[the formalised variant]\label{rem:leanH}
The Lean formulation (\S\ref{sec:lean}) replaces (H1) and the product clause of (H2) by the
single unconditional congruence $S(ap+r,bp+s)\equiv S(a,b)S(r,s)\pmod p$ for all $a,b$ and
all $r,s<p$, together with $S(n,k)=0$ for $k>n$. In the carrying regime both sides vanish
mod $p$, so the dichotomy is absorbed; the double sum then factors because a \emph{full}
block $[0,(a+1)p)$ is split and both marginal sums are complete. The formalised theorem is
therefore immune to the defect of Remark~\ref{rem:H2} --- there is no product-region
argument in it at all --- at the cost of demanding the digit congruence for every $b$, not
only for $b\le a$. This is a genuine simplification and we recommend it as the definitive
form of the hypothesis.
\end{remark}

\begin{observation}[integrality]\label{obs:integrality}
\VerifiedR{$0$ failures, $15$ families $\times$ $16$ primes $\times$ $n\le400$}
$\vp(B(n))\ge-w\lfloor\log_pn\rfloor$, i.e.\ $d_n^{\,w}B(n)\in\Zp$.
\end{observation}

\begin{remark}\label{rem:integrality}
This is \emph{not} a corollary of Theorem~\ref{thm:LB}. That theorem is a single-digit
statement --- (H4) is used in the form $x(n,k)\le n<p^2$, which is precisely $a<p$ --- and
the iteration over base-$p$ digits would need it at $a\ge p$, where the descended letters
$\KK r\chi{\lfloor x/p\rfloor}$ are no longer $p$-integral. Even granting
Conjecture~\ref{conj:master}, passing to the ratio form
$p^{w}B(n)/A(n)\equiv\chi(p)B(q)/A(q)$ requires $A(n)$ and $A(q)$ to be $p$-units, which is
a separate matter (see \cite{MalikStraub} for the primes that never divide a given
sporadic sequence). Apart from the $\mathbf B$ case proved in
Theorem~\ref{thm:Bden} below, we therefore record
Observation~\ref{obs:integrality} as computed, and
in \S\ref{sec:BV} treat the implication from \eqref{eq:master} to mirror-map integrality as
a heuristic and not as a proof.
\end{remark}

\begin{theorem}[sharp denominator for $\mathbf B$]\label{thm:Bden}
For $h\in\Z$, let $B_h(0)=0$, $B_h(1)=1$, and
\[
 (n+1)^2B_h(n+1)
 =h(3n^2+3n+1)B_h(n)-3h^2n^2B_h(n-1).
\]
Then
\[
 \bigl(\forall n\ge0,\ d_n^2B_h(n)\in\Z\bigr)
 \quad\Longleftrightarrow\quad 3\mid h.
\]
At $h=3$ this is Zagier's sporadic pair $\mathbf B$.  Hence the $\mathbf B$ row of
Observation~\ref{obs:integrality} is a theorem, and its exponent $2$ is optimal.
\end{theorem}

\begin{proof}
Set
\[
 T_h(n)=\sum_{k=0}^n\binom nk(-h)^{n-k}B_h(k),
 \qquad T_h(m)=0\quad(m\le0).
\]
Conjugating the recurrence by this binomial transform gives the step-three endpoint
recurrence
\[
 n^2T_h(n)+h^3(n-2)(n-1)T_h(n-3)=(-h)^{n-1}.
\]
Indeed, at the generating-function level the original operator
\[
 \theta^2-hz(3\theta^2+3\theta+1)+3h^2z^2(\theta+1)^2
\]
is conjugated to
\[
 \theta^2+h^3z^3(\theta+1)(\theta+2).
\]
Iteration yields
\[
 (-1)^{n-1}T_h(n)=h^{n-1}
 \sum_{\substack{0\le j<n\\n-j\equiv1\ (3)}}
 \frac{\prod_{t=0}^{r-1}(j+2+3t)(j+3+3t)}
      {\prod_{t=0}^{r}(j+1+3t)^2},
 \qquad r=\frac{n-j-1}{3}.
\tag{B-end}\label{eq:Bend}
\]
All propagation signs agree because $j+r\equiv n-1\pmod2$.

Let $R(n,j)$ denote the quotient in \eqref{eq:Bend}.  If $p\ne3$, counting multiples of
each $p^a$ in the three arithmetic progressions shows
\[
 v_p(R(n,j))\ge-2v_p(d_n):
\]
the doubled denominator progression has length $r+1$, the two numerator progressions
have length $r$, and the deficit for each $p^a$ is at most $2$.  For $p=3$, choose a
factor of maximal $3$-adic valuation in $D=\prod_{t=0}^r(j+1+3t)$.  The difference of any
other factor from it is $3$ times an index difference, whence
\[
 v_3(D)\le v_3(d_n)+r+v_3(r!),
 \qquad
 v_3(R(n,j))\ge-2v_3(d_n)-2r-2v_3(r!).
\]
If $3\mid h$, then $n-1=j+3r\ge3r$ and $v_3(r!)\le r/2$, so every summand of
$d_n^2T_h(n)$ in \eqref{eq:Bend} is integral.  Binomial inversion and $d_k\mid d_n$
transfer the result to $B_h(n)$.

Conversely, direct recurrence propagation gives
\[
 B_h(6)=-\frac{39521h^5}{32400},\qquad
 d_6^2B_h(6)=-\frac{39521h^5}{9}.
\]
Since $3\nmid39521$, integrality at $n=6$ forces $3\mid h$.  Finally
$B_3(2)=21/4$, so $d_2B_3(2)=21/2\notin\Z$; the exponent $2$ is sharp.
\end{proof}

%======================================================================
\section{Status of the fifteen}\label{sec:status}

Table~\ref{tab:status} is the summary; the four subsections that follow justify it.

\begin{table}[t]
\footnotesize
\setlength{\tabcolsep}{4pt}
\begin{tabular}{@{}llccl@{}}
\toprule
family & $\mathfrak w$ known? & tame? & $\Delta\equiv0$? & status of \eqref{eq:LB}\\
\midrule
$\gamma$   & yes, $2$ terms  & yes & free & \Thm\ $p\ge5$; \Lean\\
$\mathbf A$& yes, $3$ terms  & yes & free & \Thm\ $p\ge5$ ($p\ne2$ suffices); \Lean\\
$\mathbf D$& yes, $4$ terms  & yes & free & \Thm\ $p\ge7$; $p=5$ \Verified\ only\\
$s_{10}$   & yes, $3$ terms  & yes & free & \Thm\ $p\ge7$\\[2pt]
$s_7$      & yes, $8$ terms  & no  & \Verified\ $p\le37$ & reachable; two local bounds, \Open\\[2pt]
$\alpha$   & yes, $14$ terms & no  & \Excluded & blocked; needs (LD)\\
$\varepsilon$& yes, $10$ terms& no & \Excluded & blocked; needs (LD)\\
$\mathbf E$& yes, $5$ terms, $\chi_{-4}$ & no & \Excluded\ (inv.) & blocked; needs (LD)\\[2pt]
$\mathbf B,\mathbf C,\mathbf F,\delta,\zeta,\eta,s_{18}$
           & no (\Excluded) & --- & --- & conjectural\\
\bottomrule
\end{tabular}
\medskip
\caption{Status of the second-row congruence \eqref{eq:LB}. ``Tame'' means every letter
argument lies in $[0,n]$ on the summation range, i.e.\ (H4). ``$\Delta\equiv0$ free'' means
Lemma~\ref{lem:Deltatame} applies.}
\label{tab:status}
\end{table}

\subsection{The harmonic decompositions}\label{sec:decomps}
All of the following were obtained by exact linear fitting of an ansatz
$B(n)=\sum_kS(n,k)\mathfrak w(n,k)$ with $\mathfrak w$ as in \eqref{eq:weightshape}, solved
over a large prime field, rationally reconstructed, then \emph{validated exactly over $\Q$
on held-out $n$}. Each row below is \VerifiedR{held out $n=60,\dots,72$, exact; and
re-validated for this paper at $n\le20$}.

\begin{center}\small
\begin{tabular}{@{}llll@{}}
\toprule
family & $w$ & $\mathfrak w(n,k)$ & arguments\\
\midrule
$\gamma$&3&$\tfrac13\HH3n-\tfrac16\HH3k$&$n,k$\\
$\mathbf A$&2&$\tfrac14\HH2k+\tfrac34\HH1k(\HH1k-\HH1{n-k})$&$k,n-k$\\
$\mathbf D$&2&$\tfrac15\bigl[\HH2n+\HH1k(2\HH1k-\HH1{n-k}-\HH1n)\bigr]$&$n,k,n-k$\\
$s_{10}$&2&$\tfrac15\HH2k+\tfrac45\HH1k(\HH1k-\HH1{n-k})$&$k,n-k$\\
$s_7$&2&$8$ terms in $H^{(1)},H^{(2)}$&$k,n-k,n,2k$\\
$\varepsilon$&3&$10$ terms in $H^{(1)},H^{(2)},H^{(3)}$&$k,n-k,2k,2k-n$\\
$\alpha$&3&$14$ terms in $H^{(1)},H^{(2)},H^{(3)}$&$k,n-k,2k,2n-2k$\\
$\mathbf E$&2&$\tfrac12\KK2{\chi_{-4}}{2k}
  +\bigl(\tfrac34\HH1k-\tfrac12\HH1{2k}\bigr)
   \bigl(\KK1{\chi_{-4}}{2k}-\KK1{\chi_{-4}}{2n-2k}\bigr)$&$k,2k,2n-2k$\\
\bottomrule
\end{tabular}
\end{center}

\begin{remark}[the dichotomy: a heuristic, not a theorem]\label{obs:dichotomy}
A pure harmonic monomial $\prod_t\HH{r_t}{x_t}$ produces, in the limit, values in the
$\Q$-algebra generated by $\zeta(2),\zeta(3),\dots$ and $\log2$ (the last from letters at
wide arguments, e.g.\ $\HH1{2k}-\HH1k\to\log2$). \emph{If} $L(\chi,w)$ lies outside that
algebra --- which is expected but open, already for Catalan's constant
$G=L(\chi_{-4},2)$, whose irrationality is not known --- then no $\chi$-twisted family can
have a pure-$\zeta$ harmonic decomposition, and its letters must be twisted. We label this
as what it is: an expectation resting on unavailable transcendence input. It is not used in
the proof of any theorem below; it is the explanation the data invites.
\end{remark}

This is confirmed by the decisive experiment on $\mathbf E$: \VerifiedR{bases of up to $44$
elements over $8$ argument forms} $\mathbf E$'s fitting system is \emph{inconsistent in
every pure harmonic basis tried}, and becomes consistent the moment $\chi_{-4}$-letters are
admitted, yielding the five-term weight above. $\mathbf E$ is the only known instance of the
twisted case $e=1$ of Theorem~\ref{thm:LB}.

\subsection{The four theorems}\label{sec:proved}
\begin{theorem}\label{thm:four}
\Thm\ Congruence \eqref{eq:LB} holds for $\gamma$ and $\mathbf A$ with $\chi\equiv1$ and
all $p\ge5$, and for $\mathbf D$ and $s_{10}$ with $\chi\equiv1$ and all $p\ge7$.
\end{theorem}

\begin{proof}
Apply Theorem~\ref{thm:LB} with the weights of \S\ref{sec:decomps}. (H5) holds with $e=0$.
(H4): the arguments are among $n,k,n-k$, all in $[0,n]$; the coefficients have denominators
$3,6$ ($\gamma$), $4$ ($\mathbf A$), $5$ ($\mathbf D$), $5$ ($s_{10}$), whence the stated
prime bounds. (H2)'s $\Delta$-clause is free by Lemma~\ref{lem:Deltatame}. For (H1)--(H3),
write $S=\bin nk^{\alpha}\bin{n+k}n^{\beta}$: the surviving set is $\{s\le r\}$ (no borrow
in $\bin nk$) intersected with $\{r+s<p\}$ (no carry in $\bin{n+k}n$; this second
constraint is vacuous for $\mathbf A$ and $s_{10}$, where $\beta=0$), a product region;
there $\bin nk\equiv\bin ab\bin rs$ and $\bin{n+k}n\equiv\bin{a+b}a\bin{r+s}r$ by Lucas, and
the deleted terms with $r+s\ge p$ vanish mod $p$, so
$\sum_{s\in\Sigma_r}S(r,s)\equiv A(r)$. Digit compatibility is
$\lfloor k/p\rfloor=b$, $\lfloor(n-k)/p\rfloor=a-b$ (no borrow), $\lfloor n/p\rfloor=a$.
The details are \cite[\S3]{CompanionLucas2nd}.
\end{proof}

\VerifiedR{$0$ failures, all cells $n=ap+r<p^2$} both layers of the proof separately:
$\gamma$, $\mathbf A$, $s_{10}$ at $p=5,7,11,13$; $\mathbf D$ at $p=7,11,13$, with $11$
failing cells at $p=5$ --- exactly the (H4) coefficient condition, $\mathbf D$'s
coefficients being $\tfrac15(1,2,-1,-1)$. So $\mathbf D$'s proof is unconditional for
$p\ge7$ and $p=5$ is \Verified\ only. ($s_{10}$'s coefficients are likewise
$\tfrac15(1,4,-4)$; it happens to show no $p=5$ defect, but the proof as written still
requires $p\ge7$.) \emph{No other defect of any kind appears in the tame cases: the only
obstruction the mechanism ever meets is the coefficient condition.}

\subsection{$s_7$: reachable}\label{sec:s7}
For $s_7$, $\varepsilon$, $\alpha$ and $\mathbf E$ the known weight is \emph{not} tame:
arguments reach $2n$ (letters at $2k$, $2n-2k$, $2k-n$), so for $k$ in the vanishing layer
$\lfloor x/p\rfloor$ can reach $2p-1$, the descended letter acquires a pole, and
$p^{w}\mathfrak w(n,k)$ is no longer $p$-integral. The two-layer split of the proof of
Theorem~\ref{thm:LB} then fails \emph{layer by layer}: each of
\[
  V(n)=\!\!\sum_{k:\,p\mid S(n,k)}\!\!S(n,k)p^{w}\mathfrak w(n,k),
  \qquad
  U(n)=\!\!\sum_{k:\,p\nmid S(n,k)}\!\!S(n,k)p^{w}\mathfrak w(n,k)
\]
separately misses its intended value. (It does \emph{not} fail index by index in $k$; that
misreading led to a false claim in an earlier memo of the programme, corrected in
\cite{LemmaDCheck}.)

\begin{proposition}\label{prop:s7}
\Verified\ For $s_7$ and every prime $p\in\{5,11,13,17,19,23,29,31,37,41\}$, all cells
$n=ap+r<p^2$ and all summation indices $k$:
\[
  \textup{(V1)}\quad \vp(S(n,k))+\vp\bigl(p^{2}\mathfrak w(n,k)\bigr)\ge1
  \quad\text{whenever }p\mid S(n,k),
\]
with minimum exactly $1$ ($1\,659\,068$ vanishing-layer terms, $0$ violations); and
\[
  \textup{(V2)}\quad \vp(S(a,b))+\vp(\mathfrak w(a,b))\ge1
  \quad\text{whenever }p\mid S(a,b),\ a<p,
\]
$0$ violations for $p\in\{5,11,13,17,19,23,29,31,37\}$.
\end{proposition}

\noindent
\textup{(V1)} implies that the vanishing layer is $\equiv0\pmod p$ and \textup{(V2)} that
$\Delta(a)\equiv0\pmod p$, so at each prime at which \textup{(V1)}, \textup{(V2)}, (H1),
(H3) and the product structure are known --- all \VerifiedR{$5\le p\le23$, $0$ failures}
--- the proof of Theorem~\ref{thm:LB} goes through for $s_7$. At those primes
\eqref{eq:LB} was in any case verified directly, so what \emph{is} gained is not a fact but
a route: a proof of \textup{(V1)} and \textup{(V2)} for all $p\ge5$, $p\ne7$ would make
$s_7$ a theorem. We state this as a problem and not as a proposition.

The prime $p=7$ is excluded for $s_7$: its coefficients have denominator $28$, and there are
$13$ failing cells at $p=7$ against none at $p=5,11,13$. The reason (V1) is true for $s_7$
and false for the other three is structural: in $s_7$'s weight the argument $2k$ occurs
\emph{only inside weight-one letters}, so the pole has order at most $1$ and a single Kummer
carry in $\bin{2k}n$ beats it. This is a bounded, purely local Kummer-versus-order-one-pole
problem.

\begin{problem}\label{prob:s7}
\Open\ Prove (V1) and (V2) for $s_7$, $p\ge5$, $p\ne7$. This would take the proved count
from four to five. It does \emph{not} deliver an instance with $\chi\ne1$: $s_7$ has
$\chi=\mathbf 1$.
\end{problem}

\subsection{$\alpha$, $\varepsilon$, $\mathbf E$: blocked, and provably so}\label{sec:blocked}

\begin{proposition}[the defect is real]\label{prop:defectreal}
\Verified\ For $\alpha$, $\varepsilon$ and $\mathbf E$, the inequality
$\vp(S(n,k))+\vp(p^{w}\mathfrak w(n,k))\ge1$ on $\{p\mid S(n,k)\}$ is \emph{false}, with
explicit counterexamples, and so is its aggregated weakening $\vp(V(n))\ge1$.
\end{proposition}

The minimal certificates, in exact rational arithmetic, are:
\begin{itemize}[leftmargin=1.6em]
\item $\alpha$ (Domb), $p=5$, $n=20=4\cdot5+0$, $k=5$:
  $S=\bin{20}5^2\bin{10}5\bin{30}{15}$ has $v_5(S)=1$, while
  $v_5(\mathfrak w(20,5))=-5$, so $v_5(p^3\mathfrak w)=-2$ and the sum is $-1<1$. The pole
  sits at $2n-2k=30\ge p^2$, through the monomials $\HH1k\bigl(\HH1{2n-2k}\bigr)^2$ and
  $\HH2k\HH1{2n-2k}$; the term $S\,p^{w}\mathfrak w$ is not even $p$-integral.
\item $\varepsilon$, $p=5$, $n=15$, $k=15$: $v_5(S)=2$, $v_5(p^3\mathfrak w)=-2$, sum $0<1$;
  pole at $2k=30$, through $\HH3{2k}$ and $\HH2{2k}\HH1{2k}$.
\item $\mathbf E$, $p=5$, $n=15$, $k=15$: $v_5(S)=1$, $v_5(p^2\mathfrak w)=-1$, sum $0<1$;
  pole at $2k=30$, through $\KK2{\chi_{-4}}{2k}$ and $\HH1{2k}\KK1{\chi_{-4}}{2k}$.
\end{itemize}
The observed minima are uniform in $n$ and $p$ ($-1$ for $\alpha$, $0$ for $\varepsilon$
and $\mathbf E$, \VerifiedR{$5\le p\le23$}), so a \emph{corrected} inequality with a
constant amount of slack does exist --- and is useless, because the proof needs
$\ge1$ exactly. Note the sharpening that $\varepsilon$'s summand $\bin nk^2\bin{2k}n^2$ is a
square, so $\vp(S)\in2\Z$ and one Kummer carry already buys two orders; that is why
$\varepsilon$ lands at $0$ rather than $-2$.

\begin{proposition}[what is true instead]\label{prop:LD}
\Verified\ For all four non-tame families, all cells $n=ap+r<p^2$ and
$p\in\{5,7,11,13,17,19,23\}$ ($\alpha,\varepsilon,\mathbf E$) resp.\
$p\in\{5,11,\dots,37\}$ ($s_7$), $0$ failures in $20/20$, $42/42$, $110/110$, $156/156$,
$272/272$, $342/342$, $506/506$ cells for \emph{both} lines:
\begin{align}
  V(n)&\equiv\chi(p)^{e}\,\Delta(a)\,A(r) \pmod p, \tag{LD}\label{eq:LD}\\
  U(n)&\equiv\chi(p)^{e}\,\bigl(B(a)-\Delta(a)\bigr)A(r) \pmod p. \tag{LD$'$}
\end{align}
Adding the two lines returns \eqref{eq:LB}. So \eqref{eq:LD} is exactly the missing
ingredient, and nothing else is missing.
\end{proposition}

This is the precise content of the earlier, vaguer statement that ``the vanishing-layer and
surviving-layer defects occur in equal numbers and cancel''. They are equal in number
because they are the \emph{same} number, $\chi(p)^{e}\Delta(a)A(r)$, with opposite signs.
We stress that ``equal numbers'' was never independent evidence: given the (verified)
conclusion \eqref{eq:LB}, $V$ fails if and only if $U$ does. The genuine new information is
the closed form. Certificates: $\alpha$ at $p=5$, $n=20$: $V\equiv\Delta(4)A(0)=37259/54$,
matched at $v_5=3$; $\mathbf E$ at $p=7$, $n=28$: $\Delta(4)=379/9$, $\chi_{-4}(7)=-1$,
$V\equiv-379/9$, matched at $v_7=2$. Note also that $\Delta(a)=0$ identically for
$a\le(p-1)/2$ --- the defect is a large-digit phenomenon.

\begin{theorem}[no admissible weight repairs the defect]\label{thm:norepair}
\Excluded\ Let $\mathcal K$ be the exact $\Q$-kernel of the fitting system for one of
$\alpha$, $\varepsilon$, $\mathbf E$ over the argument alphabets below, so that the
admissible weights form the affine space $\mathfrak w_0+\mathcal K$. Then, \emph{for each
of the primes $p\in\{5,7,11,13,17,19,23\}$ marked} $\times$ \emph{below}, the $\mathbb
F_p$-linear system $\sum_{b\in\mathcal B_a}S(a,b)\mathfrak w(a,b)\equiv B(a)$
($a=1,\dots,p-1$), i.e.\ $\Delta(a)\equiv0$ for all $a$, is \emph{inconsistent}:
\begin{center}\small
\begin{tabular}{@{}llccccccccc@{}}
\toprule
family & alphabet & cols & $\dim\mathcal K$ & $5$&$7$&$11$&$13$&$17$&$19$&$23$\\
\midrule
$\alpha$&$\{k,n-k,2k,2n-2k\}$&40&21&$\times$&$\times$&$\times$&$\times$&$\times$&$\times$&$\times$\\
$\alpha$&$+\,n$&65&32&\checkmark&\checkmark&$\times$&$\times$&$\times$&$\times$&$\times$\\
$\varepsilon$&$\{k,n-k,2k,2k-n\}$&40&2&$\times$&$\times$&$\times$&$\times$&$\times$&$\times$&$\times$\\
$\varepsilon$&$+\,n,\ 2n-2k$&98&7&$\times$&$\times$&$\times$&$\times$&$\times$&$\times$&$\times$\\
$\mathbf E$&$\{k,n-k,2k,2n-2k\}+\chi_{-4}$, (H5)&20&10&$\times$&$\times$&$\times$&$\times$&$\times$&$\times$&$\times$\\
$\mathbf E$&$+\,n$, (H5) dropped&65&28&$\times$&$\times$&$\times$&$\times$&$\times$&$\times$&$\times$\\
$s_7$&$\{n,k,n-k,2k\}$ ($+$enlarged)&14/27&2&\checkmark&$\times$&\checkmark&\checkmark&\checkmark&\checkmark&\checkmark\\
\bottomrule
\end{tabular}
\end{center}
For $\mathbf E$ the kernel has rank $0$ against $\Delta$ at each of those primes:
\emph{every} admissible weight produces the \emph{same} $\Delta(a)\bmod p$.
\end{theorem}

\begin{remark}
The obvious extrapolation --- that $\Delta(a)\bmod p$ is an invariant of $\mathbf E$ at
every prime, not an artefact of the fit --- is not proved here and is recorded as such.
Nothing below depends on it: what the theorem already excludes is the strategy of
re-choosing the weight, and it excludes it at every prime that has been examined.
\end{remark}

\begin{proof}[Method and scope]
The kernel basis is computed exactly (row reduction modulo four primes, CRT, rational
reconstruction) and each basis vector is verified over $\Q$ by
$\sum_kS(n,k)\kappa(n,k)=0$ for $n=1,\dots,21$ ($0$ failures in every case); the
$\mathbb F_p$ system is then solved exactly. The $\mathbf E$ statement is confirmed
independently by direct perturbation: at $p=5$ all ten kernel vectors leave
$(\Delta(3),\Delta(4))\equiv(3,1)\bmod 5$ unchanged; for $\alpha$ only two kernel vectors
move it, and only along $(1,-1)$, never onto $(0,0)$.
\emph{Scope, stated honestly:} this excludes every weight-homogeneous $\Q$-combination of
bare letters at the listed argument forms. It does not exclude a wholly different ansatz ---
rational-function coefficients, a constant-term representation, or new letters.
\end{proof}

\begin{corollary}\label{cor:count}
The route ``prove one valuation inequality and move $\alpha,\varepsilon,s_7,\mathbf E$
inside Theorem~\ref{thm:LB}'' --- proposed in an earlier memo of the programme and repeated
in \cite[\S9]{CompanionLucas2nd} --- is wrong. The inequality is false for
$\alpha,\varepsilon,\mathbf E$; it works for $s_7$ alone, taking the proved count from four
to five, not from four to eight. For $\alpha$, $\varepsilon$, $\mathbf E$ the target is
\eqref{eq:LD}, which is a \emph{descent} congruence of the same logical type as the theorem
it is meant to prove, not a local Kummer estimate. That is a materially larger gap, and it
is where the programme's first instance with $\chi\ne1$ is gated.
\end{corollary}

Two structural leads survive. First, $\Delta(a)=0$ identically for $a\le(p-1)/2$. Second,
the per-digit refinement $V_b(n)\equiv\chi(p)^{e}S(a,b)\mathfrak w(a,b)A(r)$ of
\eqref{eq:LD} is \VerifiedR{$0$ failures, $p=5,7,11,13$} for $\varepsilon$, $\mathbf E$ and
$s_7$, so for those (LD) is provable one digit at a time; it is \emph{false} for $\alpha$
(e.g.\ $16$ of $30$ defect-digits fail at $p=5$), so $\alpha$ needs the aggregated form.

Finally, $\mathbf E$ is provably non-tame in a stronger sense: attempts to find a
\emph{tame} $\chi_{-4}$-decomposition for $\mathbf E$ over the complete tame argument set
$\{n,k,n-k\}$, with $\chi_{-4}$, with three characters, with and without (H5), are all
\Excluded\ at $106$--$154$ excess equations; the only other tame representation of
$\mathbf E$'s first row that exists, $A(n)=\sum_k4^{\,n-2k}\bin n{2k}\bin{2k}k^2$
(\VerifiedR{$n\le30$}), supports no decomposition at all ($230$ columns, rank $230$, $90$
excess); and a sweep of $4.4$ million tame templates finds no third representation. The wide
arguments are essential, so the twisted case of Theorem~\ref{thm:LB} is proved as a
mechanism but has, at present, no tame instance.

\subsection{The conjectural seven}\label{sec:seven}
For $\mathbf B$, $\mathbf C$, $\mathbf F$, $\delta$, $\zeta$, $\eta$, $s_{18}$ no harmonic
decomposition is known and \eqref{eq:master} is \Verified\ only. The searches were run to
completion in every \emph{complete tame} alphabet (the set of tame integer linear forms on a
given support is finite and explicitly computable), at up to $520$ columns and with $55$ to
$193$ excess equations; all $30$ runs are inconsistent, and each headline negative was
recomputed at a second, unrelated modulus with identical rank and verdict. Two of these are
worth stating as results.

\begin{proposition}[$\mathbf C$]\label{prop:C}
\Excluded\ $\mathbf C$'s second solution admits no harmonic-monomial decomposition on the
summand $\bin nk^2\bin{2k}k$ in any basis tried: plain harmonic monomials at $\{n,k,n-k,2k\}$
($14$ elements) and at $\{n,k,n-k,2k,n+k,2n-2k,2n,2n-k\}$ ($44$), and with
$\chi_{-3}$-letters added ($44$ and $152$ elements, $110$ and $220$ exact equations); and,
sharper, not even when the conductor-three arguments $3k$, $3n-3k$, $3n$ that the summand
does not supply are furnished by hand ($90$ columns / $64$ excess and $132$ / $63$). By
the heuristic of Remark~\ref{obs:dichotomy} its letters are expected to be twisted; the
obstruction is therefore not merely the missing arguments. Moreover $\mathbf C$'s complete tame form set is only
$\{n,k,n-k\}$ (its $2k$ reaches $2n$), so Theorem~\ref{thm:LB} could not be applied to
$\mathbf C$ even if a decomposition were found in a wide alphabet.
\end{proposition}

\begin{proposition}[$\delta$]\label{prop:delta}
\Excluded\ $\delta$'s second solution admits no weight-three harmonic decomposition over its
complete tame alphabet $\{n,k,2k,3k,n-k,n-2k,n-3k\}$ at full degree ($140$ columns, $60$
excess; also $330$ columns / $70$ excess with $\chi_{-3}$-letters), and $\delta$ possesses no
tame binomial representation whatsoever in a $4.4$-million-template sweep.
\end{proposition}

A useful diagnostic emerged from these runs: the rank deficiency of the design matrix.
Families that do have a decomposition have large kernels ($\mathbf E$: rank $37$ of $65$;
$\alpha$: $19$ of $40$); $\mathbf C$, $\mathbf F$, $s_{18}$, $\eta$ have small ones
($86/90$, $100/152$, $107/110$, $153/156$); and $\mathbf B$, $\delta$, $\zeta$ have
\emph{rank $=$ columns in every run} --- no $\Q$-relation at all between their summand and
the harmonic letters. So enlarging the alphabet will not help those three; the ansatz has to
change. The recommended replacement is Gorodetsky's constant-term representation
\cite{Gorodetsky}, which supplies exactly the sort of new letters that
Proposition~\ref{prop:C} shows are needed.

The open coverage gaps, stated rather than hidden: degree-three-in-weight-one coverage for
$\zeta$ and $\eta$ in their complete tame spaces, and for $\mathbf F$ at seven tame
arguments, would need $424$--$1088$ columns and correspondingly $500$--$1150$ rows; those
three runs at full degree were not made.

%======================================================================
\section{Closed forms}\label{sec:closed}

\subsection{Ap\'ery's $\zeta(3)$ numerators}
\begin{theorem}\label{thm:gamma}
\Thm\ \Lean\ Let $b_n$ be defined by $b_0=0$, $b_1=6$ and
$(n+2)^3b_{n+2}=(34n^3+153n^2+231n+117)\,b_{n+1}-(n+1)^3b_n$, i.e.\ Ap\'ery's recurrence.
Then for every $n\ge0$
\[
  b_n=\sum_{k=0}^n\bin nk^2\bin{n+k}k^2\bigl(2\HH3n-\HH3k\bigr).
\]
\end{theorem}

The proof is in \cite{CompanionLucasMin}; it is elementary --- two Zeilberger/Gosper
certificates that are polynomial identities in $\Q[n,k]$, combined into a \emph{single}
telescope with vanishing boundary values --- and it is formalised (\S\ref{sec:lean}). It
replaces both Ap\'ery's own weight
$\HH3n+\sum_{m\le k}(-1)^{m-1}\bigl(2m^3\bin nm\bin{n+m}m\bigr)^{-1}$ and Nesterenko's
weight $\HH3k+(2\HH1k-\HH1{n+k}-\HH1{n-k})\HH2k$ \cite{Nesterenko1996}: the fitting system
has an $11$-dimensional kernel, so the harmonic-monomial representation is far from unique,
and this is the shortest member. We have not located it in the literature; see
\cite[\S9]{CompanionLucasMin} for the search record.

\subsection{The symmetric template}\label{sec:template}
For the $d$-fold sums $A^{(d)}(n)=\sum_k\bin nk^d$ the summand is symmetric under
$k\mapsto n-k$, so any fitted weight may be replaced by its $k\leftrightarrow n-k$ average
without changing $B(n)$. Writing
\begin{equation}\label{eq:sigmadelta}
  \sigma(n,k)=\HH2k+\HH2{n-k},\qquad \delta(n,k)=\HH1k-\HH1{n-k},
\end{equation}
the fitted weights of \S\ref{sec:decomps} symmetrise to
$\mathfrak w=\alpha_d\sigma+\beta_d\delta^2$ with
$\alpha_d=\tfrac1{2(d+1)}$, $\beta_d=\tfrac d{2(d+1)}$; note $\alpha_d+\beta_d=\tfrac12$ for
every $d$.

This is not cosmetic. In the symmetric form the only weight-one letter that ever appears is
the single antisymmetric combination $\delta$, because $\sigma$ differences to
\emph{rational} functions under $n\mapsto n\pm1$ and $k\mapsto k+1$:
\begin{equation}\label{eq:diffs}
\begin{aligned}
  \sigma(n+1,k)-\sigma(n,k)&=\frac{1}{(n+1-k)^2},&
  \delta(n+1,k)-\delta(n,k)&=\frac{-1}{n+1-k},\\
  \sigma(n-1,k)-\sigma(n,k)&=\frac{-1}{(n-k)^2},&
  \delta(n-1,k)-\delta(n,k)&=\frac{1}{n-k},\\
  \sigma(n,k+1)-\sigma(n,k)&=\tfrac1{(k+1)^2}-\tfrac1{(n-k)^2},&
  \delta(n,k+1)-\delta(n,k)&=e(n,k):=\tfrac1{k+1}+\tfrac1{n-k}.
\end{aligned}
\end{equation}
So the certificate module is free of rank $2$ on $\{1,\delta\}$ instead of rank $4$ on
$\{1,\HH1k,\HH1{n-k},\HH1n\}$: \emph{two} Gosper problems instead of four.

\medskip\noindent\textbf{The template.} Let $c_+(n)u_{n+1}+c_0(n)u_n+c_-(n)u_{n-1}=0$ be the
recurrence and $S(n,k)=\bin nk^d$.
\begin{enumerate}[leftmargin=2em,label=(\arabic*)]
\item \emph{Zeilberger certificate.} Seek $\rho\in\Q(n)[k]$ with
  $G(n,k)=S(n,k)k^d\rho(n,k)/(n+1-k)^d$ satisfying
  \begin{equation}\label{eq:Z}
    c_+S(n+1,k)+c_0S(n,k)+c_-S(n-1,k)=G(n,k+1)-G(n,k).
  \end{equation}
  Because $S(n,k+1)/S(n,k)=\bigl((n-k)/(k+1)\bigr)^d$, one has
  $G(n,k+1)=S(n,k)\rho(n,k+1)$: \emph{the shifted certificate is the summand times a
  polynomial}. This is what makes the scheme work, and it turns \eqref{eq:Z} into one
  identity in $\Q(n,k)$.
\item \emph{Split and Abel.} With $\Delta^{\pm}$ the $n$-differences \eqref{eq:diffs},
  \[
    L[B](n)=\sum_{k=0}^{n}\mathfrak w(n,k)\bigl(G(n,k+1)-G(n,k)\bigr)
            +\sum_{k}\mathrm{Corr}(n,k)+\mathrm{Top}(n),
  \]
  $\mathrm{Corr}=c_+S(n+1,k)\Delta^+\mathfrak w+c_-S(n-1,k)\Delta^-\mathfrak w$ and
  $\mathrm{Top}(n)=c_+(n)S(n+1,n+1)\mathfrak w(n+1,n+1)$; Abel summation and $G(n,0)=0$
  turn the first sum into
  $\mathfrak w(n,n)G(n,n+1)-\sum_{k=0}^{n-1}G(n,k+1)\Delta_k\mathfrak w$.
\item \emph{Two Gosper problems.} By \eqref{eq:diffs} the residual $k$-summand is
  $S(n,k)\bigl(\mathcal A(n,k)+\mathcal B(n,k)\delta(n,k)\bigr)$ with
  $\mathcal A,\mathcal B\in\Q(n,k)$. Seek $Z(n,k)=S(n,k)(z_0+z_1\delta)$ with
  $Z(n,k+1)-Z(n,k)$ equal to it. Since $\delta(n,k+1)=\delta(n,k)+e(n,k)$, matching
  coefficients of $1$ and $\delta$ gives, with $r=S(n,k+1)/S(n,k)$,
  \begin{align}
    r(n,k)\,z_1(n,k+1)-z_1(n,k)&=\mathcal B(n,k),\label{eq:G1}\\
    r(n,k)\,z_0(n,k+1)-z_0(n,k)&=\mathcal A(n,k)-r(n,k)z_1(n,k+1)e(n,k).\label{eq:G2}
  \end{align}
  The ansatz $z_1=N(k)/(n+1-k)^{d+1}$, $z_0=M(k)/(n+1-k)^{d+2}$ reduces both to linear
  systems over $\Q(n)$.
\item \emph{Boundary.} Everything cancels iff $Z(n,0)=0$ and three scalar identities hold.
\end{enumerate}

\begin{proposition}[the boundary identities, for every $d$]\label{prop:boundary}
\Proved\ With $\mathfrak w=\alpha\sigma+\beta\delta^2$, and provided \eqref{eq:Z},
\eqref{eq:G1}, \eqref{eq:G2} and $Z(n,0)=0$ hold, one has $L[B](n)=0$ for all $n\ge1$
as soon as
\begin{equation}\label{eq:boundary}
\begin{gathered}
  \rho(n,n+1)=-c_+(n),\qquad
  N(n)=2\beta\,c_+(n)\Bigl((n+1)^d-\tfrac1{n+1}\Bigr),\\
  M(n)=-c_+(n)(\alpha+\beta)\Bigl((n+1)^d+\tfrac1{(n+1)^2}\Bigr).
\end{gathered}
\end{equation}
In particular, for $c_+(n)=(n+1)^{d-1}$, $\alpha=\alpha_d$, $\beta=\beta_d$,
\[
  \rho(n,n+1)=-(n+1)^{d-1},\qquad
  N(n)=\tfrac{d}{d+1}\bigl((n+1)^{2d-1}-(n+1)^{d-2}\bigr),
\]
\[
  M(n)=-\tfrac12\bigl((n+1)^{2d-1}+(n+1)^{d-3}\bigr).
\]
\end{proposition}

\begin{proof}
Assemble step (2)--(4). Since $S(n-1,k)=0$ for $k\ge n$ and $S(n,n)=S(n+1,n+1)=1$,
\[
  L[B](n)=Z(n,n)-Z(n,0)+\mathrm{Corr}(n,n)+\mathfrak w(n,n)G(n,n+1)+\mathrm{Top}(n).
\]
Now $\sigma(n,n)=\HH2n$, $\delta(n,n)=\HH1n$, $G(n,n+1)=\rho(n,n+1)$,
$Z(n,n)=M(n)+N(n)\HH1n$, $\mathrm{Corr}(n,n)=c_+(n+1)^d\bigl[\alpha+\beta(1-2\HH1n)\bigr]$
and $\mathrm{Top}(n)=c_+\bigl[\alpha\HH2{n+1}+\beta(\HH1{n+1})^2\bigr]$. Substituting
$\HH2{n+1}=\HH2n+(n+1)^{-2}$ and $\HH1{n+1}=\HH1n+(n+1)^{-1}$ and collecting the
coefficients of $\HH2n$, $(\HH1n)^2$, $\HH1n$ and $1$ shows that \eqref{eq:boundary} makes
each of the four coefficients vanish identically --- the first two both reduce to
$\rho(n,n+1)+c_+(n)=0$ --- whence $L[B](n)=0$. The last statement follows by substituting
$\alpha_d+\beta_d=\tfrac12$.
\end{proof}

\begin{remark}
The conditions \eqref{eq:boundary} are also necessary, since $\HH2n$, $(\HH1n)^2$, $\HH1n$
and $1$ are linearly independent over $\Q(n)$ --- a standard fact of $\Pi\Sigma$-difference
field theory, cf.\ \cite{Schneider2007} --- but we do not use necessity and do not prove it
here.
\end{remark}

\begin{remark}
Proposition~\ref{prop:boundary} says the boundary is \emph{never} the obstruction: the three
identities hold for every $d$. The only $d$-dependent step is the existence of a rational
solution of \eqref{eq:G1}--\eqref{eq:G2}. It is worth noting where this leaves the general
$d$: Chamberland--Straub's conjecture \cite[Conj.~9]{ChamStraub} was proved for every $d$ by
Straub--Zudilin \cite{StraubZudilin} by an entirely different route (Remark~\ref{rem:SZ}),
so the template is not needed for it; what the template supplies, where it closes, is a
\emph{finite certificate} valid for every $n\ge1$.
\end{remark}

\subsection{Franel and $\sum_k\bin nk^4$: finite certificates}\label{sec:franel}

The two closed forms below are not new --- see Remark~\ref{rem:SZ} --- but the proofs are,
and it is the harmonic-monomial \emph{display} that \S\ref{sec:status} needs.

\begin{theorem}[Franel]\label{thm:franel}
\Thm\ Let $A(n)=\sum_k\bin nk^3$ (A000172), so $(n+1)^2u_{n+1}=(7n^2+7n+2)u_n+8n^2u_{n-1}$.
Then the second solution is
\[
  B(n)=\sum_{k=0}^n\bin nk^3\Bigl[\tfrac18\bigl(\HH2k+\HH2{n-k}\bigr)
        +\tfrac38\bigl(\HH1k-\HH1{n-k}\bigr)^2\Bigr].
\]
\end{theorem}

\begin{theorem}[$s_{10}$]\label{thm:s10}
\Thm\ Let $A(n)=\sum_k\bin nk^4$, so
$(n+1)^3u_{n+1}=(2n+1)(6n^2+6n+2)u_n+n(64n^2-4)u_{n-1}$. Then the second solution is
\[
  B(n)=\sum_{k=0}^n\bin nk^4\Bigl[\tfrac1{10}\bigl(\HH2k+\HH2{n-k}\bigr)
        +\tfrac25\bigl(\HH1k-\HH1{n-k}\bigr)^2\Bigr].
\]
\end{theorem}

\begin{proof}[Proof of Theorems \ref{thm:franel} and \ref{thm:s10}]
Both sides vanish at $n=0$ and equal $1$ at $n=1$, so it suffices to prove that the right-hand
side is annihilated by the recurrence operator for $n\ge1$; the leading coefficient
$(n+1)^{d-1}$ never vanishes, so the second solution is then determined. Apply the template
of \S\ref{sec:template} with the certificates
\begin{align*}
 \rho_F(n,k)&=-\tfrac1n\bigl(4-12k+12k^2-4k^3+22n-39kn+18k^2n+32n^2-27kn^2+14n^3\bigr),\\
 N_F(k)&=\tfrac{3k^2}{4n}\bigl(4-16k+24k^2-16k^3+4k^4+26n-68kn+63k^2n-20k^3n\\
    &\qquad\qquad+54n^2-88kn^2+39k^2n^2+46n^3-36kn^3+14n^4\bigr),\\
 M_F(k)&=-\tfrac{k}{2n}\bigl(2-10k+20k^2-20k^3+10k^4-2k^5+15n-50kn+70k^2n-45k^3n\\
    &\qquad\quad+11k^4n+40n^2-90kn^2+80k^2n^2-25k^3n^2+50n^3-70kn^3+30k^2n^3\\
    &\qquad\quad+30n^4-20kn^4+7n^5\bigr)
\end{align*}
for $d=3$, and the analogous $\rho_{10}$ (degree $5$ in $k$, degree $6$ in $n$), $N_{10}$,
$M_{10}$ (degree $9$ in $k$) for $d=4$, recorded in the data file accompanying
\cite{MinimalFormMemo}. Every one of \eqref{eq:Z}, \eqref{eq:G1}, \eqref{eq:G2} is an identity
in $\Q(n,k)$ and was verified symbolically, by clearing denominators and expanding to $0$;
$z_1(n,0)=z_0(n,0)=0$ because $N$ and $M$ carry factors $k^2$, $k$ ($d=3$) resp.\ $k^3$,
$k^2$ ($d=4$). The three boundary identities \eqref{eq:boundary} hold:
$\rho_F(n,n+1)=-(n+1)^2$, $N_F(n)=\tfrac34\bigl((n+1)^5-(n+1)\bigr)$,
$M_F(n)=-\tfrac12\bigl((n+1)^5+1\bigr)$; and $\rho_{10}(n,n+1)=-(n+1)^3$,
$N_{10}(n)=\tfrac45\bigl((n+1)^7-(n+1)^2\bigr)$,
$M_{10}(n)=-\tfrac12\bigl((n+1)^7+(n+1)\bigr)$ --- exactly the values forced by
Proposition~\ref{prop:boundary}. Proposition~\ref{prop:boundary} then gives $L[B]=0$.
\end{proof}

\begin{remark}[certificate audit]\label{rem:audit}
The identities \eqref{eq:Z}, \eqref{eq:G1}, \eqref{eq:G2}, $Z(n,\cdot)$ at $k=0$, and
\eqref{eq:boundary} were each re-derived and re-verified symbolically for the present paper,
independently of the memo they come from, in a computer-algebra system, each returning $0$
after simplification. In addition $L[B]=0$ was checked in exact rational arithmetic for
$n\le15$ for both families, and $B(0)=0$, $B(1)=1$ by hand.
\end{remark}

\begin{remark}[these sums are Straub--Zudilin's $t^2$-coefficients]\label{rem:SZ}
Straub--Zudilin \cite{StraubZudilin} study the analytic deformation
\[
  A^{(s)}(n,t)=\sum_{k=0}^n\bin nk^{s}
    \prod_{j=1}^{k}\Bigl(1-\frac tj\Bigr)^{-s}\prod_{j=1}^{n-k}\Bigl(1+\frac tj\Bigr)^{-s}
  \;=\;\sum_{j\ge0}A^{(s)}_j(n)\,t^{2j},
\]
and prove (their Theorem 1.3) that every $A^{(s)}_j(n)$ solves the telescoping recurrence of
$A^{(s)}(n)$ for large enough $n$, with $A^{(s)}_1(1)=s(s+1)$. Since
\[
  \prod_{j\le k}\Bigl(1-\tfrac tj\Bigr)^{-s}\prod_{j\le n-k}\Bigl(1+\tfrac tj\Bigr)^{-s}
  =\exp\Bigl(s\,t\,\delta+\tfrac{s}{2}t^{2}\sigma+O(t^{3})\Bigr)
  =1+st\delta+t^{2}\Bigl(\tfrac s2\sigma+\tfrac{s^{2}}2\delta^{2}\Bigr)+O(t^{3})
\]
with $\sigma,\delta$ as in \eqref{eq:sigmadelta}, one gets
$A^{(s)}_1(n)/\bigl(s(s+1)\bigr)=\sum_k\bin nk^{s}\bigl(\alpha_s\sigma+\beta_s\delta^2\bigr)$
identically. \emph{So Theorems~\ref{thm:franel} and~\ref{thm:s10} are the cases $s=3,4$ of
\cite[Thm.~1.3]{StraubZudilin}, written out in harmonic numbers}
(\VerifiedR{exact, $s=3,4,5,6$, $n\le12$}). What the proof above adds is that the identity
holds for \emph{every} $n\ge1$ --- \cite{StraubZudilin} obtains the recurrence only for $n$
large --- by a finite certificate rather than a locally uniform analytic limit; that
certificate is what makes the statement formalisable in the sense of \S\ref{sec:lean}.
\end{remark}

\begin{corollary}[Ap\'ery limits, after Straub--Zudilin]\label{cor:limits}
\Proved\ For $d=3$ and $d=4$,
\[
  \lim_{n\to\infty}\frac{B^{(d)}(n)}{A^{(d)}(n)}=\frac{\zeta(2)}{d+1},
\]
i.e.\ $\pi^2/24$ and $\pi^2/30$.
\end{corollary}

\begin{proof}
Write $m=\lfloor n/2\rfloor$ and $T=\lceil n^{3/4}\rceil$. By Hoeffding's bound
$\bin nk\le2^ne^{-2(k-n/2)^2/n}$ and $\bin nm\ge2^n/(n+1)$, so
$\bin nk\le(n+1)\bin nme^{-2(k-n/2)^2/n}$ for all $k$. Hence, with
$A^{(d)}(n)\ge\bin nm^d$,
\[
  \frac1{A^{(d)}(n)}\sum_{|k-n/2|\ge T}\bin nk^d
  \;\le\;(n+1)^d\!\!\sum_{|u|\ge T}\!e^{-2du^2/n}
  \;\le\;\frac{(n+1)^d\,n}{d\,(T-1)}\,e^{-2d(T-1)^2/n}
  \;=\;O\bigl(e^{-\sqrt n}\bigr),
\]
the sum over $u\in\tfrac12\Z$ being bounded by twice the corresponding integral. On the
complementary range $|k-n/2|<T$ we have, for $n$ large, $k\ge n/3$ and $n-k\ge n/3$, whence
$|\HH2k-\zeta(2)|\le1/k\le3/n$ and likewise for $n-k$, so $|\sigma(n,k)-2\zeta(2)|\le6/n$;
and, assuming $k\ge n-k$ without loss of generality,
$|\delta(n,k)|=\sum_{j=n-k+1}^{k}j^{-1}\le(2k-n)/(n-k)\le 6n^{-1/4}$, so
$\delta(n,k)^2\le36n^{-1/2}$. Globally $0\le\sigma\le2\zeta(2)$ and
$|\delta|\le\HH1n\le1+\log n$. Therefore
\[
  \Bigl|\frac{B^{(d)}(n)}{A^{(d)}(n)}-2\alpha_d\zeta(2)\Bigr|
  \;\le\;\frac{6\alpha_d}{n}+36\beta_d n^{-1/2}
       +O\bigl(e^{-\sqrt n}\bigr)\bigl(4\alpha_d\zeta(2)+\beta_d(1+\log n)^2\bigr)
  \;\longrightarrow\;0 ,
\]
and $2\alpha_d\zeta(2)=\zeta(2)/(d+1)$. Theorems~\ref{thm:franel} and~\ref{thm:s10} identify
the sums with the second solutions normalised as in Definition~\ref{def:B}, which is the
normalisation of \cite[eq.~(20)]{ChamStraub}.
\end{proof}

\begin{remark}
Both limits were observed numerically by Cusick and are reported by van der Poorten
\cite[p.~202]{vdPoorten1979}; the case $d=3$ was proved by Zagier \cite{Zagier2009} using
modular forms, and \emph{both cases, indeed all $d$}, are Example 1.5 of Straub--Zudilin
\cite{StraubZudilin}, which settles \cite[Conj.~9]{ChamStraub}. Corollary~\ref{cor:limits}
is therefore a re-proof, not a new result. It is nevertheless worth recording, for two
reasons: it is elementary, using no modular forms and no analytic deformation, only the
closed form plus a Chernoff estimate; and it exhibits the mechanism --- the summand
$\bin nk^d$ concentrates at $k=n/2$, where $\delta(n,k)\to0$ and $\sigma(n,k)\to2\zeta(2)$,
so that the Ap\'ery limit is simply $2\alpha_d\zeta(2)$ and the $\delta^2$-part contributes
nothing. The argument does not reach $d\ge5$: the concentration estimate is uniform in $d$,
but the minimal recurrence then has order $>2$ and the closed form of
\S\ref{sec:template} is not available. \emph{Equation and conjecture numbers quoted from
\cite{ChamStraub} are those of its arXiv version; the journal version may renumber.}
\end{remark}

\subsection{A provable insufficiency: the $\zeta(2)$-Ap\'ery pair}\label{sec:Dobstruction}
\begin{observation}\label{obs:Dform}
\VerifiedR{$n\le25$} For $\mathbf D$, $A(n)=\sum_k\bin nk^2\bin{n+k}k$ (A005258),
$(n+1)^2u_{n+1}=(11n^2+11n+3)u_n+n^2u_{n-1}$, the second solution is
\[
  B(n)=\tfrac15\sum_k\bin nk^2\bin{n+k}k
        \Bigl[\HH2n+\HH1k\bigl(2\HH1k-\HH1{n-k}-\HH1n\bigr)\Bigr],
\]
with $B(0),\dots,B(6)=0,1,\tfrac{25}4,\tfrac{1741}{36},\tfrac{6585}{16},
\tfrac{13327519}{3600},\tfrac{124308457}{3600}$.
\end{observation}

Steps (1)--(3) of the template go through, and beautifully. The Zeilberger certificate is
the smallest in the programme,
\[
  \rho_D(n,k)=k^2+k(1+6n)-4-15n-11n^2,\qquad
  G_D(n,k)=\frac{k^3\rho_D(n,k)\bin{n+1}k^2\bin{n+k-1}k}{n(n+1)^2},
\]
\Proved\ (a rational identity, expanded to $0$; re-verified here). The structural boundary
identity holds \emph{exactly}:
$\rho_D(n,n+1)=-2(n+1)(2n+1)=-(n+1)^2\bin{2n+2}{n+1}\big/\bin{2n}n=-c_+(n)S(n+1,n+1)/S(n,n)$.
The correction decomposes as $Y(n,k)/S(n,k)=\mathcal A+\mathcal B_1\HH1k
+\mathcal B_2(\HH1{n-k}+\HH1n)$ with $\mathcal B_2=\rho_D(n,k+1)/(5(k+1))$
\VerifiedR{exact, $n=2,\dots,8$}.

\begin{theorem}[the $\HH1{n-k}$ component is not Gosper-summable]\label{thm:gosper}
\Proved\ Regard $n$ as an indeterminate and put $T(k)=S_D(n,k)\mathcal B_2(n,k)$. Then
$T(k+1)/T(k)=\bigl(a(k)/b(k)\bigr)\bigl(c(k+1)/c(k)\bigr)$ with
$a=(n-k)^2(n+k+1)$, $b=(k+1)^2(k+2)$, $c(k)=\rho_D(n,k+1)$, and
$\gcd\bigl(a(k),b(k+j)\bigr)=1$ in $\Q(n)[k]$ for every $j\ge0$ (the roots of $a$ are $n$
and $-n-1$, those of $b(k+j)$ are $-j-1$ and $-j-2$), so this is Gosper's normal form.
Gosper's key equation $a(k)x(k+1)-b(k-1)x(k)=c(k)$ has no solution $x\in\Q(n)[k]$.
Consequently $T$ has no hypergeometric antidifference with certificate rational in $n$,
which is what a proof uniform in $n$ requires.
\end{theorem}

\begin{proof}
Write $A=a(k)$, $B=b(k-1)$ and use
\[
  Ax(k+1)-Bx(k)=\tfrac12(A-B)\bigl(x(k+1)+x(k)\bigr)+\tfrac12(A+B)\bigl(x(k+1)-x(k)\bigr).
\]
Here
\[
  A-B=-n\bigl(k^2+(n+2)k-n(n+1)\bigr),\qquad
  A+B=2k^3+(2-n)k^2-(n^2+2n)k+(n^3+n^2),
\]
so $\deg(A-B)=2<3=\deg(A+B)$ and one is in the second case of Gosper's degree bound. For
$x$ of degree $m$ with leading coefficient $x_m$, the coefficient of $k^{m+2}$ on the left
is $\tfrac12\bigl((-n)\cdot2x_m+2\cdot mx_m\bigr)=x_m(m-n)$, which is nonzero in $\Q(n)$
because $n$ is an indeterminate. Hence $\deg\bigl(Ax(k+1)-Bx(k)\bigr)=m+2$, forcing
$m=\deg c-2=0$. With $x=x_0$ constant, the coefficient of $k^2$ in $x_0(A-B)=c(k)$ gives
$x_0=-1/n$; the coefficient of $k$ is then $n+2$, against the required $6n+3$, and
$n+2\ne6n+3$ in $\Q(n)$. Contradiction. \Verified: solving the linear system directly also
returns no solution for any $\deg x\le8$.
\end{proof}

\begin{remark}
The restriction to $\Q(n)$ is not cosmetic: for a \emph{fixed} positive integer $n$ the
degeneracy $m=n$ is not excluded and $\gcd\bigl(a(k),b(k+n)\bigr)\ne1$. The theorem rules
out a certificate valid uniformly in $n$ --- the only kind that can prove a closed form ---
and says nothing about individual values of $n$.
\end{remark}

\begin{remark}[diagnosis]
In the Franel and $s_{10}$ cases the weight is $k\leftrightarrow n-k$ symmetric, the
certificate module has rank $2$ on $\{1,\delta\}$, and each component telescopes
separately. For $\mathbf D$ the weight is genuinely three-lettered ($\HH1k$, $\HH1{n-k}$,
$\HH1n$) and the components do not telescope separately, while the weight-two part of any
enlarged ansatz is forced to vanish ($\Delta_k(Sz)=0\Rightarrow z=0$). Closing $\mathbf D$
therefore needs either a different gauge for $\mathfrak w$ --- the harmonic-monomial fit has
a large kernel, and a kernel shift changes the correction --- or a genuine $\Pi\Sigma$-field
ansatz allowing nested sums, as in Schneider's treatment of Ap\'ery's double sum
\cite{Schneider2007}. \Open. This is the first place in the programme where the rank-two
mechanism is \emph{provably} insufficient, and the obstruction is exactly the one met at
weight $5$: two independent letters at the top.
\end{remark}

\subsection{$\mathbf E$: a twisted closed form}\label{sec:E}
\begin{observation}\label{obs:Eform}
\VerifiedR{$n\le50$ by direct re-implementation; $n\le20$ re-checked here; held out
$n=70,\dots,82$ in the original fit}
For $\mathbf E$, $A(n)=\sum_k\bin nk\bin{2k}k\bin{2(n-k)}{n-k}$, the second solution is
\[
  B(n)=\sum_k\bin nk\bin{2k}k\bin{2(n-k)}{n-k}\,
       \Bigl[\tfrac12\KK2{\chi_{-4}}{2k}
       +\bigl(\tfrac34\HH1k-\tfrac12\HH1{2k}\bigr)
        \bigl(\KK1{\chi_{-4}}{2k}-\KK1{\chi_{-4}}{2n-2k}\bigr)\Bigr],
\]
with $B(0),B(1),B(2),B(3)=0,1,7,\tfrac{404}9$, and $B(n)/A(n)\to G/2$, $G$ Catalan's
constant.
\end{observation}

This is the only known closed form for a sporadic second solution with a non-trivial
character, and it is the empirical anchor of the whole $\chi$-twist story: the same fitting
problem is inconsistent in every pure harmonic basis and consistent as soon as
$\chi_{-4}$-letters are admitted --- exactly what the heuristic of
Remark~\ref{obs:dichotomy} predicts for a limit of the form $L(\chi_{-4},2)/2$. Every argument of the weight is wide
($2k$, $2n-2k$), and by the exhaustion recorded in \S\ref{sec:blocked} no tame alternative
exists; this is why $\mathbf E$ sits outside Theorem~\ref{thm:LB}.

%======================================================================
\section{Relation to Beukers--Vlasenko's Conjecture 7.5}\label{sec:BV}

Set $f(x)=1-t\,g(x)$ for a Laurent polynomial $g$ and $F(t)=\sum_nf_nt^n$ with
$f_n=\ct\bigl(g(x)^n\bigr)$; for the sporadic families these $f_n$ are the $A(n)$
\cite{Gorodetsky,Straub2301}. Let $F_N(t)$ be the truncation of $F$ at $t^N$, let
$F(t)\log t+G(t)$ be the second solution of the Picard--Fuchs operator --- so that the
coefficients of $G$ are our $B(n)$, up to normalisation --- and let
$q(t)=t\exp\bigl(G(t)/F(t)\bigr)$ be the canonical (mirror) coordinate. Beukers--Vlasenko
prove in \cite[\S7]{BVthree}, building on \cite{BVtwo}, that for all $m,s\ge1$,
\begin{equation}\label{eq:BV76}
  \frac{F(t)}{F(t^\sigma)}\;\equiv\;\frac{F_{mp^s}(t)}{F_{mp^{s-1}}(t^\sigma)}\pmod{p^s},
\end{equation}
whose coefficientwise reading with the naive lift $t^\sigma=t^p$ and $s=1$ is the first-row
Lucas congruence \eqref{eq:Arow}. Their Conjecture 7.5 asserts \eqref{eq:BV76} modulo
$p^{2s}$ for the \emph{excellent} Frobenius lift --- the unique lift with
$q^{\sigma}=\gamma^{p-1}q^{p}$ --- in the completely symmetric Calabi--Yau families; they
state that they could not prove it for a single $g$.

The relationship to \eqref{eq:master} is as follows.

\begin{enumerate}[leftmargin=2em]
\item \textbf{Two rows of one matrix.} The crystal $\mathrm{CY}(g)(2)$ of \cite{BVthree} is
  free of rank $2$ on $1/f$, $\theta(1/f)$, and its Frobenius has diagonal
  $(\lambda_0,\lambda_2)$ with $\lambda_0=F(t)/F(t^\sigma)$ the unit root. Conjecture 7.5 is
  a statement about $\lambda_0$ alone, sharpened from $p^s$ to $p^{2s}$. Statement
  \eqref{eq:master} is the statement about the \emph{other} diagonal entry: that Frobenius
  acts on the second graded piece with eigenvalue $\chi(p)p^{w}$. In coefficient form,
  their $\lambda_0$-congruence \emph{is} the $A$-row and ours \emph{is} the $B$-row.
  \textbf{Neither formally implies the other.}
\item \textbf{But 7.5's standing hypothesis is a weak form of our conclusion.}
  Conjecture 7.5 is stated only for the excellent lift, which exists as a lift of
  $\Z_p[[t]]$ precisely because $q(t)=t\exp(G/F)\in\Z_p[[t]]$ --- $p$-integrality of the
  mirror map, a statement about $G$, i.e.\ about $B(n)$. Observation~\ref{obs:integrality}
  --- $\vp(B(n))\ge-w\lfloor\log_pn\rfloor$ --- is exactly the coefficient-level shadow of
  that integrality, and it is sharper: it does not merely assert $p$-integrality, it pins
  the depth at exactly $w$, and \eqref{eq:master} exhibits the eigenvalue $\chi(p)$ on top
  of it. Heuristically, then,
  \[
    \eqref{eq:master}\ \rightsquigarrow\ \text{coefficient form of mirror integrality}\
    \Longrightarrow\ \text{the excellent lift is defined over $\Z_p[[t]]$,}
  \]
  the last being the standing hypothesis of 7.5. The first arrow is written with a wavy
  line deliberately: as Remark~\ref{rem:integrality} explains, deducing
  Observation~\ref{obs:integrality} from \eqref{eq:master} needs both the master form and a
  $p$-unit hypothesis on $A(n)$, so this is a statement about where the two circles of
  ideas sit relative to one another, not a chain of implications. In that sense
  \eqref{eq:master} is \emph{upstream} of 7.5, and certainly not equivalent to it.
\item \textbf{The character is invisible in unit-root formulations.} No character can appear
  in $\lambda_0=F/F^{\sigma}$: the unit root of the holomorphic period is a unit, and the
  first-row Lucas congruence is untwisted for all fifteen families
  \cite{MalikStraub,Straub2301}. The character surfaces only on the \emph{non-unit}
  eigenvalue. Any crystal-theoretic proof of \eqref{eq:master} must therefore carry the
  information that for seven of the fifteen families the second Frobenius eigenvalue is
  $\chi(p)p^{w}$ and not $p^{w}$, i.e.\ that the family is attached to a
  quadratic-twisted modular parametrisation. \emph{This is a concrete correction to make to
  any ``Dwork crystals $\Rightarrow$ Lucas congruences for the second solution'' slogan.}
  We have not located a second-solution version of \eqref{eq:BV76} in \cite{BVthree}, nor
  elsewhere; we do not claim to have surveyed \cite{BVone,BVtwo} exhaustively.
\item \textbf{What \eqref{eq:master} does not give.} Nothing towards the modulo-$p^{2s}$
  \emph{doubling} that is the actual content of 7.5. Our floor is exactly $w$ and equality
  is attained in every cell tested, so there is no hidden supercongruence in the $B$-row of
  the kind 7.5 predicts for the $A$-row.
\end{enumerate}

Proposition~\ref{prop:nolimit} is the sharpest evidence for the Frobenius reading. For
$\mathbf B$, $\delta$ and $\eta$ the archimedean Ap\'ery limit does not exist --- the
characteristic roots are complex conjugates of equal modulus and $B(n)/A(n)$ oscillates
forever --- yet \eqref{eq:master} holds flat for all three, with $\chi=\chi_{-3}$,
$\mathbf 1$, $\chi_5$. The $p$-adic Ap\'ery limit survives where the archimedean one does
not; whatever \eqref{eq:master} is measuring, it is not the limit.

%======================================================================
\section{Machine verification}\label{sec:lean}

The following are formalised in Lean~4 with Mathlib and checked by the kernel with no
\texttt{sorry} and no \texttt{native\_decide}; for each of them
\[
  \texttt{\#print axioms}\;\longrightarrow\;
  [\texttt{propext},\ \texttt{Classical.choice},\ \texttt{Quot.sound}].
\]

\begin{center}\footnotesize
\begin{tabular}{@{}ll@{}}
\toprule
declaration & statement\\
\midrule
\texttt{theorem\_LB} & Theorem~\ref{thm:LB} in the abstract form of Remark~\ref{rem:leanH}\\
\texttt{bMin\_eq\_bApery} & $\sum_k\bin nk^2\bin{n+k}k^2(2\HH3n-\HH3k)=b_n$, $b_n$ defined by Ap\'ery's recurrence\\
\texttt{apery\_b\_harmonic\_closed\_form} & the same, with the harmonic sums written out\\
\texttt{bMin\_lucas} & $p^3\,\tilde b_{ap+r}\equiv\tilde b_a\,a_r\pmod p$ for the explicit sum, any prime $p$\\
\texttt{bApery\_lucas} & the same for the sequence defined by Ap\'ery's recurrence\\
\texttt{bFranel\_lucas} & $p^2\,\beta_{ap+r}\equiv\beta_a f_r\pmod p$, $p\ne2$\\
\texttt{apery\_lucas} & Gessel's congruence $a_{ap+r}\equiv a_aa_r\pmod p$, in digit-product form\\
\texttt{Q\_lucas} & its analogue for a Brown--Zudilin double sum (not a sporadic first row)\\
\bottomrule
\end{tabular}
\end{center}

Three scope remarks, stated because they are easy to overstate.
\begin{enumerate}[leftmargin=2em]
\item \texttt{bMin\_lucas} carries \emph{no} hypothesis $p\ge5$: the minimal weight's
  coefficients are the integers $2,-1$, so the (H4) coefficient clause is vacuous. The
  $p\ge5$ of the informal statement is an artefact of Ap\'ery's original weight, which
  carries $1/(2m^3)$. Similarly \texttt{bFranel\_lucas} needs only $p\ne2$.
\item That \texttt{bFranel} is the second solution of the Franel recurrence is
  \Verified\ upstream and \emph{not} formalised; the Lean theorem is about the explicit sum.
  Theorem~\ref{thm:franel} of the present paper supplies the missing identification on
  paper, but it is not machine-checked. The corresponding identification \emph{is}
  formalised in the $\gamma$ case (\texttt{bMin\_eq\_bApery}).
\item The equivalence of the minimal form with van der Poorten's classical double sum
  is proved on paper \cite[\S5]{MinimalFormMemo} and is \emph{not} formalised.
\end{enumerate}
Nothing in \S\ref{sec:status} beyond Theorem~\ref{thm:LB} and its two instances is
formalised; in particular Theorem~\ref{thm:norepair}, Theorem~\ref{thm:gosper} and
Corollary~\ref{cor:limits} are paper proofs.

A standing lesson of the programme belongs here: \texttt{\#print axioms} cannot detect a
false undischarged hypothesis. In a different part of the development a hypothesis that was
false as stated left everything downstream ``axiom-clean''; only an attempt to discharge it
caught the error. The house rule that followed --- no certificate reaches the kernel without
a faithful re-implementation outside Lean --- is why Remark~\ref{rem:audit} exists.

%======================================================================
\section{Open problems}\label{sec:open}

Ranked by the ratio of value to distance.

\begin{enumerate}[leftmargin=2em,label=\textbf{P\arabic*.}]
\item \textbf{$s_7$, two local bounds.} Problem~\ref{prob:s7}: prove (V1) and (V2) for
  $s_7$, $p\ge5$, $p\ne7$. Bounded, purely local, Kummer-versus-order-one-pole; both are
  \Verified\ with $0$ violations in $1.66$ million vanishing-layer terms and the minimum is
  exactly $1$, so nothing is being asked for beyond what the data already shows. Proved
  count $4\to5$.
\item \textbf{The descent congruence \eqref{eq:LD}.} Prove
  $V(n)\equiv\chi(p)^{e}\Delta(a)A(r)\pmod p$ for $\alpha$, $\varepsilon$, $\mathbf E$. This
  is where the first instance with $\chi\ne1$ is gated. Two leads: $\Delta(a)=0$ for
  $a\le(p-1)/2$, so it is a large-digit phenomenon; and the per-digit refinement is
  \Verified\ for $\varepsilon$ and $\mathbf E$ (not for $\alpha$), which makes those two
  provable one digit at a time. By Theorem~\ref{thm:norepair} re-fitting the weight is
  excluded, so the proof must use something the fit does not see.
\item \textbf{The master form.} Prove Conjecture~\ref{conj:master} --- the mod-$p^w$,
  all-$n$ statement --- in even one case. Theorem~\ref{thm:LB} gives only the single-digit
  mod-$p$ shadow. The gap is the same one left open at weight $3$ in
  \cite{CompanionLucas2nd}.
\item \textbf{New letters for the conjectural seven.} By the rank diagnostic of
  \S\ref{sec:seven}, $\mathbf B$, $\delta$, $\zeta$ admit no $\Q$-relation at all between
  summand and harmonic letters, so no enlargement of a $\sum_kS\cdot\mathfrak w$ fit will
  succeed. The two candidates are (a) Gorodetsky-style constant-term representations
  \cite{Gorodetsky}, now known to be necessary rather than optional for $\mathbf C$
  (Proposition~\ref{prop:C}), and (b) rational-function coefficients $c_j(n,k)$, which would
  require replacing (H4)'s coefficient clause by $c_j(n,k)\equiv c_j(a,b)\pmod p$ on the
  surviving set --- note that the obvious candidates $1/(2k+1)$ and $1/(n+1)$ \emph{fail}
  that test.
\item \textbf{$\mathbf D$'s closed form.} Upgrade Observation~\ref{obs:Dform} to a theorem.
  Theorem~\ref{thm:gosper} says the rank-two mechanism cannot do it; try a gauge change
  inside the $11$-dimensional kernel of the fit, or a $\Pi\Sigma$ ansatz
  \cite{Schneider2007}. This is a miniature of the weight-$5$ blockage in a case small
  enough to compute, and it should be settled before anyone spends weeks at weight $5$.
\item \textbf{Deformations as a source of letters.} Straub--Zudilin's deformation
  $A^{(s)}(n,t)$ (Remark~\ref{rem:SZ}) produces the entire harmonic weight of
  $\sum_k\bin nk^{s}$ as one Taylor coefficient. Does an analogue exist for the other
  sporadic summands? This is a strictly larger class than the one already excluded in
  \S\ref{sec:seven}: the programme's earlier ``$\Gamma$-derivation'' probe covered only
  \emph{degree-one} weights $\sum_x\lambda_x\HH wx$ and found $\gamma$ alone consistent,
  whereas the deformation above yields the degree-two weight
  $\alpha\sigma+\beta\delta^2$. Since $\mathbf B$, $\delta$ and $\zeta$ admit no
  $\Q$-relation at all between summand and harmonic letters, a deformation --- which
  changes the summand, not the alphabet --- is precisely the kind of move that could
  supply letters no fit can find.
\item \textbf{A crystal-theoretic proof.} Give the argument of \S\ref{sec:BV} its natural
  home: a Frobenius-eigenvalue statement on the second graded piece of
  $\mathrm{CY}(g)(2)$, carrying the character. The three families with no archimedean limit
  are the test cases, because for them no archimedean argument can be smuggled in.
\end{enumerate}

%======================================================================
\section*{Provenance and data}

Every statement above is traceable to one of the following memos of the programme, which
contain the scripts and the raw logs: \cite{LBWGeneral} (the fifteen pairs, the sweeps, the
decompositions, Lemma K and Theorem LB), \cite{LemmaDCheck} (the (H2) correction,
Propositions~\ref{prop:defectreal} and \ref{prop:LD}, Theorem~\ref{thm:norepair}),
\cite{SporadicBare} (the complete tame searches, Propositions~\ref{prop:C} and
\ref{prop:delta}, the rank diagnostic), \cite{MinimalFormMemo} (\S\ref{sec:closed}). All
arithmetic underlying a \Verified\ label is exact --- \texttt{Fraction} arithmetic or
fixed-precision $p$-adics cross-validated against it cell by cell --- and no verdict rests
on a single modulus: every negative in \S\ref{sec:seven} and \S\ref{sec:blocked} was
recomputed at a second, unrelated prime with identical rank and verdict. For the present
paper the following were re-derived independently: the recurrences of all fifteen binomial
sums ($n\le12$); the five decompositions of \S\ref{sec:decomps} ($n\le20$); the twisted
congruence \eqref{eq:LB} for all fifteen families at $p=5,7,11,13$ ($4920$ cells, $0$
failures); the master floor for all fifteen at $p=5,7$ and all $n<p^3$ (exactly $w$ in all
$30$ cells); the discriminants of Proposition~\ref{prop:nolimit}; the weights of
Observation~\ref{obs:weight} ($n\le120$); and every certificate identity of
\S\ref{sec:template}--\S\ref{sec:Dobstruction}.

%======================================================================
\begin{thebibliography}{99}

\bibitem{AlmkvistZudilin2006}
G.~Almkvist and W.~Zudilin,
\emph{Differential equations, mirror maps and zeta values},
in: Mirror Symmetry V, AMS/IP Stud.\ Adv.\ Math.\ \textbf{38}, Amer.\ Math.\ Soc.,
Providence, RI, 2006, 481--515; \texttt{arXiv:math/0402386}.

\bibitem{ASvSZ}
G.~Almkvist, D.~van~Straten and W.~Zudilin,
\emph{Ap\'ery limits of differential equations of order 4 and 5},
in: Modular Forms and String Duality, Fields Inst.\ Commun.\ \textbf{54}, Amer.\ Math.\
Soc., Providence, RI, 2008, 105--123.

\bibitem{BVone}
F.~Beukers and M.~Vlasenko,
\emph{Dwork crystals I},
Int.\ Math.\ Res.\ Not.\ (2021), no.~12, 8807--8844; \texttt{arXiv:1903.11155}.

\bibitem{BVtwo}
F.~Beukers and M.~Vlasenko,
\emph{Dwork crystals II},
Int.\ Math.\ Res.\ Not.\ (2021), 4427--4444; \texttt{arXiv:1907.10390}.

\bibitem{BVthree}
F.~Beukers and M.~Vlasenko,
\emph{Dwork crystals III: from excellent Frobenius lifts towards supercongruences},
Int.\ Math.\ Res.\ Not.\ (2023), doi:10.1093/imrn/rnad101; \texttt{arXiv:2105.14841}.

\bibitem{ChamStraub}
M.~Chamberland and A.~Straub,
\emph{Ap\'ery limits: experiments and proofs},
Amer.\ Math.\ Monthly \textbf{128} (2021), no.~9, 811--824; \texttt{arXiv:2011.03400}.

\bibitem{Cooper2012}
S.~Cooper,
\emph{Sporadic sequences, modular forms and new series for $1/\pi$},
Ramanujan J.\ \textbf{29} (2012), no.~1--3, 163--183.

\bibitem{Cooper2302}
S.~Cooper,
\emph{Ap\'ery-like sequences defined by four-term recurrence relations},
\texttt{arXiv:2302.00757}.

\bibitem{Gessel1982}
I.~M.~Gessel,
\emph{Some congruences for Ap\'ery numbers},
J.\ Number Theory \textbf{14} (1982), no.~3, 362--368.

\bibitem{Gorodetsky}
O.~Gorodetsky,
\emph{New representations for all sporadic Ap\'ery-like sequences, with applications to
congruences},
Exp.\ Math.\ \textbf{32} (2023), no.~4, 641--656; \texttt{arXiv:2102.11839}.

\bibitem{MalikStraub}
A.~Malik and A.~Straub,
\emph{Divisibility properties of sporadic Ap\'ery-like numbers},
\texttt{arXiv:1508.00297}.

\bibitem{Nesterenko1996}
Yu.~V.~Nesterenko,
\emph{A few remarks on $\zeta(3)$},
Mat.\ Zametki \textbf{59} (1996), no.~6, 865--880;
English transl.\ Math.\ Notes \textbf{59} (1996), 625--636.

\bibitem{OsburnSahuStraub}
R.~Osburn, B.~Sahu and A.~Straub,
\emph{Supercongruences for sporadic sequences},
Proc.\ Edinburgh Math.\ Soc.\ (2) \textbf{59} (2016), no.~2, 503--518;
\texttt{arXiv:1312.2195}.

\bibitem{Schneider2007}
C.~Schneider,
\emph{Ap\'ery's double sum is plain sailing indeed},
Electron.\ J.\ Combin.\ \textbf{14} (2007), \#N5, 3 pp.

\bibitem{Shvets2026}
A.~Shvets,
\emph{The Domb Ap\'ery-limit and a proof of the Ramanujan Machine conjecture Z2},
\texttt{arXiv:2604.06239}.

\bibitem{Straub2301}
A.~Straub,
\emph{Gessel--Lucas congruences for sporadic sequences},
Monatsh.\ Math.\ (2023), doi:10.1007/s00605-023-01894-3; \texttt{arXiv:2301.12248}.

\bibitem{StraubZudilin}
A.~Straub and W.~Zudilin,
\emph{Sums of powers of binomials, their Ap\'ery limits, and Franel's suspicions},
Int.\ Math.\ Res.\ Not.\ \textbf{2023}, no.~11, 9861--9879, doi:10.1093/imrn/rnac125;
\texttt{arXiv:2112.09576}.

\bibitem{vdPoorten1979}
A.~van der Poorten,
\emph{A proof that Euler missed\dots\ Ap\'ery's proof of the irrationality of $\zeta(3)$},
Math.\ Intelligencer \textbf{1} (1979), no.~4, 195--203.

\bibitem{Yang2008}
Y.~Yang,
\emph{Ap\'ery limits and special values of $L$-functions},
J.\ Math.\ Anal.\ Appl.\ (2008), doi:10.1016/j.jmaa.2008.01.094; \texttt{arXiv:0709.1968}.

\bibitem{Zagier2009}
D.~Zagier,
\emph{Integral solutions of Ap\'ery-like recurrence equations},
in: J.~Harnad and P.~Winternitz (eds.), Groups and Symmetries: From Neolithic Scots to John
McKay, CRM Proc.\ Lecture Notes \textbf{47}, Amer.\ Math.\ Soc., 2009, 349--366.

\medskip
\hrule
\medskip

\bibitem{CompanionLucas2nd}
[Author],
\emph{$\chi$-twisted Lucas congruences for the second solutions of Ap\'ery-like
recurrences}, companion preprint.

\bibitem{CompanionLucasMin}
[Author],
\emph{A closed form and a Lucas congruence for Ap\'ery's companion sequence},
companion preprint.

\bibitem{LBWGeneral}
[Author], memo \texttt{LBW\_GENERAL} (2026): the fifteen pairs, the exact sweeps, the
harmonic decompositions, Lemma K and Theorem LB, with scripts and logs.

\bibitem{LemmaDCheck}
[Author], memo \texttt{LEMMAD\_CHECK} (2026): the (H2) correction, the counterexample
certificates, the cancellation identity (LD), the exact-kernel no-repair computation.

\bibitem{SporadicBare}
[Author], memo \texttt{SPORADIC\_BARE} (2026): the complete tame alphabets, the thirty
exclusion runs, the tame-representation sweep, the rank diagnostic.

\bibitem{MinimalFormMemo}
[Author], memo \texttt{MINIMAL\_FORM\_PROOF} (2026): the closed-form proofs of
\S\ref{sec:closed} with all certificates, and the literature ledger.

\end{thebibliography}

\end{document}
